\section{AIME I 2013}
        \begin{enumerate}
        	\item (\textit{AIME I 2013}) The AIME Triathlon consists of a half-mile swim, a 30-mile bicycle ride, and an eight-mile run. Tom swims, bicycles, and runs at constant rates. He runs fives times as fast as he swims, and he bicycles twice as fast as he runs. Tom completes the AIME Triathlon in four and a quarter hours. How many minutes does he spend bicycling?


 
 \textbf{Answer: }150



\textbf{Solution}

Let $r$ represent the rate Tom swims in miles per minute. Then we have

 $\frac{1/2}{r} + \frac{8}{5r} + \frac{30}{10r} = 255$

 Solving for $r$, we find $r = 1/50$, so the time Tom spends biking is $\frac{30}{(10)(1/50)} = \boxed{150}$ minutes.

 




	\item (\textit{AIME I 2013}) Find the number of five-digit positive integers, $n$, that satisfy the following conditions:


 <ul>
(a) the number $n$ is divisible by $5,$
</ul><ul>
(b) the first and last digits of $n$ are equal, and
</ul><ul>
(c) the sum of the digits of $n$ is divisible by $5.$
</ul>
 \textbf{Answer: }200



\textbf{Solution 1}

The number takes a form of $\overline{5xyz5}$, in which $5|(x+y+z)$. Let $x$ and $y$ be arbitrary digits. For each pair of $x,y$, there are exactly two values of $z$ that satisfy the condition of $5|(x+y+z)$. Therefore, the answer is $10\times10\times2=\boxed{200}$

 





\textbf{Solution 2 (casework)}

We know the number will take the form $\overline{5xyz5}$ because of the first two conditions. The third condition means that $5|(x+y+z)$, where $x,y,z$ are nonnegative integers less than $10$. Let's split the problem into cases, where each case represents a possible sum of $x,y,z$. 

 1. If $x+y+z=0$, we only have $x=0,y=0,z=0$, so only $1$.

 2. If $x+y+z=5$, we use stars & bars. We have $5$ stars and $3-1=2$ bars, so this case has $7\choose{2}$$= 21$.

 3. If $x+y+z=10$, we use similar logic. We have $10$ stars and $2$ bars, so $12\choose{2}$$= 66$. However, $x,y,z$ must be less than $10$. Three of our order pairs have $10: (10,0,0), (0,10,0), (0,0,10)$. Therefore, this case has $66-3=63$.

 4. If $x+y+z=15$, it gets more complicated. Using the same system as used previously would be too complicated. But remember that this case is equivalentto if we wanted to choose $x+y+z=12$. We know this because you can take any ordered pair satisfying $x+y+z=15,$ subtract each variable from $9$, and get an ordered pair satisfying $x+y+z=12$. For example, take $(5,6,4)$, which satisfies $x+y+z=15.$ Its corresponding ordered pair would be $(4,3,5)$, which satisfies $x+y+z=12$. Let's proceed to calculating. Applying stars and bars, we get $14\choose{2}$$= 91$, but we have to subtract the subcases including $10,11,$ or $12$ because $x,y,z$ must all be one-digit integers. There are $3$ cases with a $12$, $6$ cases with an $11$, and $3+6$ cases with a $10$. So this case has $91-18=73$.

 5. If $x+y+z=20$, we can just calculate the ways to get $x+y+z=7$. Applying stars & bars, we get $9\choose{2}$$= 36$.

 6. If $x+y+z=25$, we can just calculate the ways to get $x+y+z=2.$ Applying stars & bars, we get $4\choose{2}$$= 6$.

 Therefore, our answer is $1+21+63+73+36+6 = \boxed{200}$.

 Note: For case 4, the subcases that must be excluded are $(12,0,0), (11,1,0), (10,2,0), (10,1,1)$ and each of their respective permutations. Those $4$ ordered pairs have $3,6,6,3$ permutations respectively, which is why $18$ ordered pairs must be subtracted from $91$.

 





\textbf{Solution 3 (less casework)}

Again, from the first two conditions we know the number takes the form $\overline{5xyz5}$, and the third condition implies that $5|(x+y+z)$. Assume that $x$, $y$, and $z$ are less than 5. Now, we will proceed with casework on the sum $x+y+z$. 

 Case 1: $x+y+z=0$. Since $x$, $y$, and $z$ are nonnegative, there is only $1$.

 Case 2: $x+y+z=5$. Using stars and bars, we obtain $\binom{5+3-1}{3-1} = 21$. However, we assumed $x, y, z < 5$, so we subract 3 to get 18.

 Case 3: $x+y+z=10$. Our only possible triplets are $(4, 4, 2)$ and $(4, 3, 3)$ and their rearrangements, totalling $3 \cdot 2 = 6$.

 Due to our assumption that $x, y, z < 5$, $x+y+z$ cannot be 15, so our casework is done. Summing, we find a total of $1+18+6=25$. However, for each of $x$, $y$, and $z$, we can add $5$ without violating the conditions. So, we can either choose to add $5$ or not, so we have to multiply $2$ to our cases three times for $x$, $y$, and $z$. This yields $25 \cdot 2 \cdot 2 \cdot 2 = \boxed{200}$.

 




	\item (\textit{AIME I 2013}) Let $ABCD$ be a square, and let $E$ and $F$ be points on $\overline{AB}$ and $\overline{BC},$ respectively. The line through $E$ parallel to $\overline{BC}$ and the line through $F$ parallel to $\overline{AB}$ divide $ABCD$ into two squares and two nonsquare rectangles. The sum of the areas of the two squares is $\frac{9}{10}$ of the area of square $ABCD.$ Find $\frac{AE}{EB} + \frac{EB}{AE}.$


 



 
 \textbf{Answer: }018



\textbf{Solution}

It's important to note that $\dfrac{AE}{EB} + \dfrac{EB}{AE}$ is equivalent to $\dfrac{AE^2 + EB^2}{(AE)(EB)}$

 We define $a$ as the length of the side of larger inner square, which is also $EB$, $b$ as the length of the side of the smaller inner square which is also $AE$, and $s$ as the side length of $ABCD$. Since we are given that the sum of the areas of the two squares is$\frac{9}{10}$ of the the area of ABCD, we can represent that as $a^2 + b^2 = \frac{9s^2}{10}$. The sum of the two nonsquare rectangles can then be represented as $2ab  = \frac{s^2}{10}$. 

 Looking back at what we need to find, we can represent $\dfrac{AE^2 + EB^2}{(AE)(EB)}$ as $\dfrac{a^2 + b^2}{ab}$. We have the numerator, and dividing$\frac{s^2}{10}$ by two gives us the denominator $\frac{s^2}{20}$. Dividing $\dfrac{\frac{9s^2}{10}}{\frac{s^2}{20}}$ gives us an answer of $\boxed{018}$.

 





\textbf{Solution 2}

Let the side of the square be $1$.  Therefore the area of the square is also $1$.   We label $AE$ as $a$ and $EB$ as $b$.  Notice that what we need to find is equivalent to: $\frac{a^2+b^2}{ab}$.Since the sum of the two squares ($a^2+b^2$) is $\frac{9}{10}$ (as stated in the problem) the area of the whole square, it is clear that the sum of the two rectangles is $1-\frac{9}{10} \implies \frac{1}{10}$.  Since these two rectangles are congruent, theyeach have area:  $\frac{1}{20}$.  Also note that the area of this is $ab$.  Plugging this into our equation we get:

 $\frac{\frac{9}{10}}{\frac{1}{20}} \implies \boxed{018}$

 


 





\textbf{Solution 3}

Let $AE$ be $x$, and $EB$ be $1$. Then we are looking for the value $x+\frac{1}{x}$. The areas of the smaller squares add up to $9/10$ of the area of the large square, $(x+1)^2$. Cross multiplying and simplifying we get $x^2-18x+1=0$. Rearranging, we get $x+\frac{1}{x}=\boxed{018}$

 





\textbf{Solution 4 (Vieta's)}

As before, $\dfrac{AE}{EB} + \dfrac{EB}{AE}$ is equivalent to $\dfrac{AE^2 + EB^2}{(AE)(EB)}$. Let $x$ represent the value of $AE=CF$. Since $EB=FB=1-x,$ the area of the two rectangles is $2x(1-x)=-2x^2+2x=\frac1{10}$. Adding $2x^2-2x$ to both sides and dividing by $2$ gives $x^2-x+\frac1{20}=0.$ Note that the two possible values of $x$ in the quadratic both sum to $1,$ like how $AE$ and $EB$ does. Therefore, $EB$ must be the other root of the quadratic that $AE$ isn't. Applying Vietas and manipulating the numerator, we get $\frac{x_1^2+x_2^2}{x_1x_2}=\frac{(x_1+x_2)^2-2x_1x_2}{\frac{1}{20}}=\frac{1^2-\frac1{10}}{\frac1{20}}=\frac{\frac9{10}}{\frac{1}{20}}=\boxed{018}$.

 





\textbf{Solution 5 (Fast)}

Let $AE = x$ and $BE = y$. From this, we get $AB = x + y$. The problem is asking for $\frac{x}{y} + \frac{y}{x}$, which can be rearranged to give $\frac{x^2 + y^2}{xy}$. The problem tells us that $x^2 + y^2 = \frac{9(x+y)^2}{10}$. We simplify to get $x^2 + y^2 = 18xy$. Finally, we divide both sides by $xy$ to get $\frac{x^2 + y^2}{xy} = \boxed{018}$. - Spacesam

 





\textbf{Solution 5 (A faster Vieta's)}

After we get the polynomial $x^2 - 18x + 1,$ we want to find $x + \frac 1 {x}.$ Since the product of the roots of the polynomial is 1, the roots of the polynomial are simply $x, \frac 1 {x}.$ Hence $x + \frac 1 {x}$ is just $18$ by Vieta's formula, or $\boxed{018}$

 





\textbf{Solution 6}

We have the equation $x^2 + y^2$ = $\frac {9}{10} \cdot (x+y)^2$. We get $x^2 + y^2 = 18xy$. We rearrange to get $x^2 + y^2 - 18xy = 0$. Since the problem only asks us for a ratio, we assume $x$ = $1$. We have $y^2 - 18y + 1$ = $0$. Solving the quadratic yields $9 + 4 \sqrt 5$ and $9 - 4 \sqrt 5$. It doesn't really matter which one it is, since both of them are positive. We will use $9 + 4 \sqrt 5$. 

 We have $9 + 4 \sqrt 5 + \frac {1}{9+4 \sqrt 5}$. Rationalizing the denominator gives us $9 + 4 \sqrt 5 + \frac {9 - 4 \sqrt 5}{81-80} = (9 + 4 \sqrt 5) + (9 - 4 \sqrt 5) = 18$. Our answer is $\boxed {018}$

 





\textbf{Solution 7}

Set side length of square to be $10$, $AE = x$ and $EB = y$. From this, we get $y+x=10$, and since the area of the square will be 100, the area of the two rectangles will be $2xy = 10$. We can substitute and say that $2xy = x+y$, and subtract $y$ from both sides, and then divide by $y$, getting the equation $\frac {x}{y} = 2x-1$, and doing the same thing with $x$ to get $\frac {y}{x} = 2y-1$. Adding these equations, we get the desired sum to be $2(x+y) - 2$, or $20-2$ which is equal to $\boxed {018}$.

 





\textbf{Solution 8 (Quadratic Equation)}

You can call the side length of the square 1. Call length AE x and solve the quadratic equation x\textasciicircum{}2+1-2x+x\textasciicircum{}2=9/10 with the quadratic formula.

 




	\item (\textit{AIME I 2013}) In the array of $13$ squares shown below, $8$ squares are colored red, and the remaining $5$ squares are colored blue. If one of all possible such colorings is chosen at random, the probability that the chosen colored array appears the same when rotated $90^{\circ}$ around the central square is $\frac{1}{n}$ , where $n$ is a positive integer. Find $n$.


 
                    \begin{figure}[H]
                        \centering
                        \includegraphics[width=0.5\textwidth]{Problem/AIME_2013_I_4_Problem_Figure_1.png}
                    \end{figure}
                    


 
 \textbf{Answer: }429



\textbf{Solution}

When the array appears the same after a 90-degree rotation, the top formation must look the same as the right formation, which looks the same as the bottom one, which looks the same as the right one. There are four of the same configuration. There are not enough red squares for these to be all red, nor are there enough blue squares for there to be more than one blue square in each three-square formation. Thus there are 2 reds and 1 blue in each, and a blue in the center. There are 3 ways to choose which of the squares in the formation will be blue, leaving the other two red.

 There are $\binom{13}{5}$ ways to have 5 blue squares in an array of 13. 

 $\frac{3}{\binom{13}{5}}$ = $\frac{1}{429}$ , so \textit{$n$} = $\boxed{429}$

 





\textbf{Solution 2}

By the Pigeonhole Principle, if the middle square isn't blue, than any 90 degree rotation won't map onto itself because one extra blue square will be mapped onto a red square. Therefore, we have to have 1 blue square in the middle and 1 blue square in each of the 4 seperate containers(think if there is two in each a 90 degree rotation wouldn't move it to the other side). There are 3 places in each of the seperate tilings for it and everything else will be red. Therefore, there are only 3 good combinations. Now the total amount is 13choose5, so we lead to the answer: $\frac{3}{\binom{13}{5}}$ = $\frac{1}{429}$ , so \textit{$n$} = $\boxed{429}$

 





\textbf{Solution 3}

We use the same idea as the other solutions, knowing the middle square must be blue, which has probability $\frac{5}{13}$, the next blue square can go anywhere; then the last three must go in the three spots they rotate in, giving us $\frac{3!}{11*10*9}$. Multiplying yields $\boxed{429}$.

 




	\item (\textit{AIME I 2013}) The real root of the equation $8x^3-3x^2-3x-1=0$ can be written in the form $\frac{\sqrt[3]{a}+\sqrt[3]{b}+1}{c}$, where $a$, $b$, and $c$ are positive integers. Find $a+b+c$.


 



 
 \textbf{Answer: }098



\textbf{Solution 1}

We note that $8x^3 - 3x^2 - 3x - 1 = 9x^3 - x^3 - 3x^2 - 3x - 1 = 9x^3 - (x + 1)^3$. Therefore, we have that $9x^3 = (x+1)^3$, so it follows that $x\sqrt[3]{9} = x+1$. Solving for $x$ yields $\frac{1}{\sqrt[3]{9}-1} = \frac{\sqrt[3]{81}+\sqrt[3]{9}+1}{8}$, so the answer is $\boxed{98}$.

 





\textbf{Solution 2}

Let $r$ be the real root of the given \href{https://artofproblemsolving.com/wiki/index.php?title=Polynomial}{polynomial}. Now define the cubic polynomial $Q(x)=-x^3-3x^2-3x+8$. Note that $1/r$ must be a root of $Q$. However we can simplify $Q$ as $Q(x)=9-(x+1)^3$, so we must have that $(\frac{1}{r}+1)^3=9$. Thus $\frac{1}{r}=\sqrt[3]{9}-1$, and $r=\frac{1}{\sqrt[3]{9}-1}$. We can then multiply the numerator and denominator of $r$ by $\sqrt[3]{81}+\sqrt[3]{9}+1$ to rationalize the denominator, and we therefore have $r=\frac{\sqrt[3]{81}+\sqrt[3]{9}+1}{8}$, and the answer is $\boxed{98}$.

 





\textbf{Solution 3}

It is clear that for the algebraic degree of $x$ to be $3$ that there exists some cubefree integer $p$ and positive integers $m,n$ such that $a = m^3p$ and $b = n^3p^2$ (it is possible that $b = n^3p$, but then the problem wouldn't ask for both an $a$ and $b$). Let $f_1$ be the \href{https://artofproblemsolving.com/wiki/index.php?title=Automorphism}{automorphism} over $\mathbb{Q}[\sqrt[3]{a}][\omega]$ which sends $\sqrt[3]{a} \to \omega \sqrt[3]{a}$ and $f_2$ which sends $\sqrt[3]{a} \to \omega^2 \sqrt[3]{a}$  (note : $\omega$ is a cubic \href{https://artofproblemsolving.com/wiki/index.php?title=Roots\_of\_unity}{root of unity}).

 Letting $r$ be the root, we clearly we have $r + f_1(r) + f_2(r) = \frac{3}{8}$ by Vieta's formulas. Thus it follows $c=8$.Now, note that $\sqrt[3]{a} + \sqrt[3]{b} + 1$ is a root of $x^3 - 3x^2 - 24x - 64 = 0$. Thus $(x-1)^3 = 27x + 63$ so $(\sqrt[3]{a} + \sqrt[3]{b})^3 = 27(\sqrt[3]{a} + \sqrt[3]{b}) + 90$. Checking the non-cubicroot dimension part, we get $a + b = 90$ so it follows that $a + b + c = \boxed{98}$.

 





\textbf{Solution 4}

We have $cx-1=\sqrt[3]{a}+\sqrt[3]{b}.$ Therefore $(cx-1)^3=(\sqrt[3]{a}+\sqrt[3]{b})^3=a+b+3\sqrt[3]{ab}(\sqrt[3]{a}+\sqrt[3]{b})=a+b+3\sqrt[3]{ab}(cx-1).$ We have 
 \[c^3x^3-3c^2x^2-(3c\sqrt[3]{ab}-3c)x-(a+b+1-3\sqrt[3]{ab})=0.\]We will find $a,b,c$ so that the equation is equivalent to the original one. Let $\dfrac{3c^2}{c^3}=\dfrac{3}{8}, \dfrac{3c\sqrt[3]{ab}-3c}{c^3}=\dfrac{3}{8}, \dfrac{a+b+1-3\sqrt[3]{ab}}{c^3}=\dfrac{1}{8}.$ Easily, $c=8, \sqrt[3]{ab}=9,$ and $a+b=90.$ So $a + b + c = 90+8=\boxed{98}$.

 




	\item (\textit{AIME I 2013}) Melinda has three empty boxes and $12$ textbooks, three of which are mathematics textbooks. One box will hold any three of her textbooks, one will hold any four of her textbooks, and one will hold any five of her textbooks. If Melinda packs her textbooks into these boxes in random order, the probability that all three mathematics textbooks end up in the same box can be written as $\frac{m}{n}$, where $m$ and $n$ are relatively prime positive integers. Find $m+n$.


 



 
 \textbf{Answer: }047



\textbf{Solution 1}

The total ways the textbooks can be arranged in the 3 boxes is $12\textbf{C}3\cdot 9\textbf{C}4$, which is equivalent to $\frac{12\cdot 11\cdot 10\cdot 9\cdot 8\cdot 7\cdot 6}{144}=12\cdot11\cdot10\cdot7\cdot3$.   If all of the math textbooks are put into the box that can hold $3$ textbooks, there are $9!/(4!\cdot 5!)=9\textbf{C}4$ ways for the other textbooks to be arranged.  If all of the math textbooks are put into the box that can hold $4$ textbooks, there are $9$ ways to choose the other book in that box, times $8\textbf{C}3$ ways for the other books to be arranged.  If all of the math textbooks are put into the box with the capability of holding $5$ textbooks, there are $9\textbf{C}2$ ways to choose the other 2 textbooks in that box, times $7\textbf{C}3$ ways to arrange the other 7 textbooks.  $9\textbf{C}4=9\cdot7\cdot2=126$, $9\cdot 8\textbf{C}3=9\cdot8\cdot7=504$, and $9\textbf{C}2\cdot 7\textbf{C}3=9\cdot7\cdot5\cdot4=1260$, so the total number of ways the math textbooks can all be placed into the same box is $126+504+1260=1890$.  So, the probability of this occurring is $\frac{(9\cdot7)(2+8+(4\cdot5))}{12\cdot11\cdot10\cdot7\cdot3}=\frac{1890}{27720}$.  If the numerator and denominator are both divided by $9\cdot7$, we have $\frac{(2+8+(4\cdot5))}{4\cdot11\cdot10}=\frac{30}{440}$.  Simplifying the numerator yields $\frac{30}{10\cdot4\cdot11}$, and dividing both numerator and denominator by $10$ results in $\frac{3}{44}$.  This fraction cannot be simplified any further, so $m=3$ and $n=44$.  Therefore, $m+n=3+44=\boxed{047}$.

 





\textbf{Solution 2}

Consider the books as either math or not-math where books in each category are indistiguishable from one another. Then, there are $\,_{12}C_{3}$ total distinguishable ways to pack the books. Now, in order to determine the desired propability, we must find the total number of ways the condition that all math books are in the same box can be satisfied. We proceed with casework for each box:

 \textbf{Case 1:}The math books are placed into the smallest box. This can be done in $\binom{3}{3}$ ways.

 \textbf{Case 2:}The math books are placed into the middle box. This can be done in $\binom{4}{3}$ ways.

 \textbf{Case 3:}The math books are placed into the largest box. This can be done in $\binom{5}{3}$ ways.

 So, the total ways the condition can be satisfied is $\binom{3}{3} + \binom{4}{3} + \binom{5}{3}$. This can be simplified to $\binom{6}{4} = \binom{6}{2}$ by the \href{https://artofproblemsolving.com/wiki/index.php?title=Hockey\_Stick\_Identity}{Hockey Stick Identity}. Therefore, the desired probability is $\dfrac{\dbinom{6}{2} }{\dbinom{12}{3}}$ = $\dfrac{3}{44}$, and $m+n=3+44=\boxed{047}$.

 





\textbf{Solution 3 (Permutation)}

There are three cases as follows. Note these are PERMUTATIONS, as the books are distinct!

 1. Math books in the 3-size box. Probability is $\frac{3\cdot2\cdot1}{12\cdot11\cdot10}$, because we choose one of the $3$ places for math book 1, then one of the $2$ for math book 2, and  the last one. Total number of orders: $12\cdot11\cdot10=1320$.

 2. In the 4-size: same logic gets you $\frac{1}{55}$, since we have $4$ places for math book 1, and so on.

 3. In the 5-size: you get $\frac{1}{22}$, for a sum of $\frac{3}{44}$ so your answer is $\boxed{047}$.

 





\textbf{Solution 4}

Assume that the $9$ other books are all distinct. The number of ways to place the other $9$ books and the $3$ is $\frac {12!}{3!}$. The number of ways to put all the $3$ math books into the box that holds $3$ books is $9!$. The number of ways to put all the $3$ math books into the box that holds $4$ books is $\binom {4}{3} \cdot 9!$, and the number of ways to put all the $3$ math books into the box that holds $5$ books is $\binom {5}{3} \cdot 9!$. The number of desired outcomes is $15 \cdot 9!$, and the total number of outcomes is $\frac {12!}{6}$. Simplifying, we get the answer is $\frac {3}{44}$, so our answer is $\boxed {047}$.

 





\textbf{Solution 5}

Assume all books are distinguishable. The total number of ways to pack the books is $\frac{12!}{3!4!5!}=27720$, where we account for symmetry of books in the same box. 

 
Case 1: All math books are in the box of capacity 3

 There are 9 books left, and the 3 boxes have 0, 4, and 5 spaces. The number of ways is $\frac{9!}{4!5!}=126$.

 
Case 2: All math books are in the box of capacity 4

 There are 9 books left, and the 3 boxes have 3, 1, and 5 spaces. The number of ways is $\frac{9!}{3!5!}=504$.

 
Case 3: All math books are in the box of capacity 5

 There are 9 books left, and the 3 boxes have 3, 4, and 2 spaces. The number of ways is $\frac{9!}{3!4!2!}=1260$.

 
Adding all these cases and dividing by the total number of ways yields $\frac{126+504+1260}{27720}=\frac{1890}{27720}=\frac{3}{44} \implies m+n=\boxed{47}$.

 


 If ur too lazy to read the above words, try the solution below!

 





\textbf{Solution 6 (No Words)}


 \[\dfrac{\frac{9!}{(3-3)!4!5!} + \frac{9!}{3!(4-3)!5!} + \frac{9!}{3!4!(5-3)!}}{\frac{12!}{3!4!5!}} = \dfrac{3}{44}\implies\boxed{47}.\]

 




	\item (\textit{AIME I 2013}) A rectangular box has width $12$ inches, length $16$ inches, and height $\frac{m}{n}$ inches, where $m$ and $n$ are relatively prime positive integers. Three faces of the box meet at a corner of the box. The center points of those three faces are the vertices of a triangle with an area of $30$ square inches. Find $m+n$.


 



 
 \textbf{Answer: }041



\textbf{Solution 1 (Pythagorean)}

Let the height of the box be $x$.

 After using the Pythagorean Theorem three times, we can quickly see that the sides of the triangle are 10, $\sqrt{\left(\frac{x}{2}\right)^2 + 64}$, and $\sqrt{\left(\frac{x}{2}\right)^2 + 36}$. Since the area of the triangle is $30$, the altitude of the triangle from the base with length $10$ is $6$.

 Considering the two triangles created by the altitude, we use the Pythagorean theorem twice to find the lengths of the two line segments that make up the base of $10$.

 We find:
 \[10 = \sqrt{\left(28+x^2/4\right)}+x/2\]

 Solving for $x$ gives us $x=\frac{36}{5}$. Since this fraction is simplified:
 \[m+n=\boxed{041}\]

 \textbf{Small Note} The 3D version of the Pythagorean theorem can also be applied here. ~MC

 





\textbf{Solution 2 (Vectors)}

We may use vectors. Let the height of the box be $2h$. Without loss of generality, let the front bottom left corner of the box be $(0,0,0)$. Let the center point of the bottom face be $P_1$, the center of the left face be $P_2$ and the center of the front face be $P_3$.

 We are given that the area of the triangle $\triangle P_1 P_2 P_3$ is $30$. Thus, by a well known formula, we note that $\frac{1}{2}|\vec{P_1P_2} \text{x} \vec{P_1P_3}|=30$ We quickly attain that $\vec{P_1P_2}=<-6,0,h>$ and $\vec{P_1P_3}=<0,-8,h>$ (We can arbitrarily assign the long and short ends due to symmetry)

 Computing the cross product, we find:
 \[\vec{P_1P_2} x \vec{P_1P_3}=-<6h,8h,48>\]

 Thus:
 \[\sqrt{(6h)^2+(8h)^2+48^2}=2*30=60\]
 \[h=3.6\]
 \[2h=7.2\]

 
 \[2h=36/5\]

 
 \[m+n=\boxed{041}\]

 





\textbf{Solution 3}

Let the height of the box be $x$.

 After using the Pythagorean Theorem three times, we can quickly see that the sides of the triangle are 10, $\sqrt{(x/2)^2 + 64}$, and $\sqrt{(x/2)^2 + 36}$. Therefore, we can use Heron's formula to set up an equation for the area of the triangle.

 The semiperimeter is $\left(10 + \sqrt{(x/2)^2 + 64} + \sqrt{(x/2)^2 + 36}\right)/2$. Therefore, when we square Heron's formula, we find

 
 \begin{align*}900 &= \frac{1}{2}\left(\left(10 + \sqrt{(x/2)^2 + 64} + \sqrt{(x/2)^2 + 36}\right)/2\right)\times\left(\left(10 + \sqrt{(x/2)^2 + 64} + \sqrt{(x/2)^2 + 36}\right)/2 - 10\right)\\&\qquad\times\left(\left(10 + \sqrt{(x/2)^2 + 64} + \sqrt{(x/2)^2 + 36}\right)/2 - \sqrt{(x/2)^2 + 64}\right)\\&\qquad\times\left(\left(10 + \sqrt{(x/2)^2 + 64} + \sqrt{(x/2)^2 + 36}\right)/2 - \sqrt{(x/2)^2 + 36}\right).\end{align*}


 Solving, we get $\boxed{041}$.

 





\textbf{Solution 4 (de Gua's)}

It isn't hard to see that the triangle connecting the centers of the faces of the rectangular prism is congruent to the triangle connecting the midpoints of three edges that concur. So we can now apply de Gua's theorem (\href{https://en.wikipedia.org/wiki/De\_Gua\%27s\_theorem}{https://en.wikipedia.org/wiki/De\_Gua\%27s\_theorem}) to see that:

 $30^2=24^2+(3x)^2+(4x)^2$

 Where $x$ is half the desired length of the height.

 Solving yields $2x=\frac{36}{5}$

 And thus $36+5=\boxed{041}$

 ---Solution 4 contributed by Siddharth Namachivayam

 





\textbf{Solution 5 (PyTRIGorean)}

Let half the height be $a$ (we want to find $2a$), then we see that the three sides of the triangle are (by Pyth Theorem) $10, \sqrt{a^2+36}, \sqrt{a^2+64}$. Using the Law of Sines with the angle as the one included between the square roots, we see that this angle's cosine is $\frac{a^2}{\sqrt{(a^2+36)(a^2+64)}}$ by the Law of Cosines, meaning that its sine is $\frac{\sqrt{100a^2+2304}}{\sqrt{(a^2+36)(a^2+64)}}$. Finally, multiply the two square-rooted sides by this sine and one-half, and equate to 30.

 You get $\sqrt{25a^2+576} = 30$, giving $a=\frac{18}{5}$. However, we must find $2a=\frac{36}{5}$ so our answer is $\boxed{041}$

 
                \begin{itemize}
                    \item Note: You can also find the cosine of the angle formed by the sides $\sqrt{a^2+64}$ and $\sqrt{a^2+36}$ using vectors. Letting the midpoint of the height of the rectangle that we are using be the origin, it follows that the 2 sides forming the angle can be represented as $<8,0,k>$ and $<0,6,k>$. Using the formula $\cos(\theta)=\frac{\mathbf{u}\cdot\mathbf{v}}{\|\mathbf{u}\|\|\mathbf{v}\|}$, we obtain $\cos(\theta)=\frac{a^2}{\sqrt{(a^2+36)(a^2+64)}}$.
                \end{itemize}


                


 





\textbf{Solution 6 (Heron's Formula Application)}

Let $x$ be $\frac12$ the height of the box. We will solve for $x$ and then multiply by $2$ at the end. The three side lengths of the triangle, by the Pythagorean Theorem, are $\sqrt{x^2+6^2},\sqrt{x^2+8^2},$ and $10$.

 Heron's formula states that the area of the triangle is $\sqrt{s(s-a)(s-b)(s-c)}$ where $s$ is the semiperimeter. Applying difference of squares to make the formula less computational, we get $\frac{\sqrt{(a+b+c)(a+b-c)(a-b+c)(-a+b+c)}}{4}=\frac{\sqrt{(a^2+b^2+2ab-c^2)(c^2-b^2-a^2+2ab)}}4=\frac{\sqrt{4a^2b^2-(a^2+b^2-c^2)^2}}4$.

 Now, we plug in $\sqrt{x^2+6^2}$ for $a$, $\sqrt{x^2+8^2}$ for $b$, and $10$ for $c$. This gives us 
 \[\frac{\sqrt{4(x^2+6^2)(x^2+8^2)-(x^2+6^2+x^2+8^2-10^2)^2}}4=30\]
 \[\sqrt{4x^4+400x^2+4\cdot48^2-4x^4}=120\]
 \[400x^2+4\cdot48^2=120^2\]
 \[25x^2+24^2=30^2\]
 \[25x^2=6^2(5^2-4^2)\]
 \[25x^2=18^2\]
 \[5x=18\]
 \[x=\frac{18}5\]

 Now multiplying by $2$, we get that the height is $\frac{36}5$, and $m+n=36+5=\boxed{041}$

 





\textbf{Solution 7 (Pyramids)}

Let the rectangular prism be $ABCDEFGH$ with $AB=12$, $AC=16$, and $AD=\alpha$ edges of the cube. It is easy to notice that the triangle given by the problem is similar to the triangle $BCD$ by a ratio of $1:2$. Thus $[BCD]=120$.

 We know that the volume of the triangular pyramid defined by $A,B,C,D$ is equal to $\frac{1}{3}\cdot12\cdot16\cdot\alpha\cdot\frac{1}{2}=32\alpha$, and we can find the volume by using the triangle $BCD$ as the base. Letting $h$ be the distance from $A$ to $\triangle BCD$, we find that the volume is $\frac{1}{3}\cdot h\cdot120=40h$. Thus $h=\frac{4}{5}\alpha$, and all that remains is to find an equation for $h$.

 Let $A$ be the origin of a three-dimensional plane, and let the line segments $AB,AC,AD$ lie on the $x,y,z$ axes respectively. We know that the equation of the plane in which $\triangle BCD$ lies can be represented by the equation $ax+by+cz=d$, where $a,b,c,d$ are real values, and we know that $B(12,0,0),C(0,16,0),D(0,0,\alpha)$ lie on $\triangle BCD$ which in turn lies on this plane. Thus the three points are solutions to the equation. Plugging in the values yields $d=12a=16b=\alpha c$, and inserting these values for $a,b,c$ in terms of $d$ and simplifying yields $\frac{1}{12}x+\frac{1}{16}y+\frac{1}{\alpha}z=1$. This is the equation of the plane.

 By the \href{https://artofproblemsolving.com/wiki/index.php?title=Distance\_formula}{formula}, the value of $h$ (or the distance between $A$ and $\triangle BCD$) is equivalent to $\frac{1}{\sqrt{\frac{1}{12^2}+\frac{1}{16^2}+\frac{1}{\alpha^2}}}$. Plugging this value in, we find that:

 
 \[\frac{4}{5}\alpha=\frac{1}{\sqrt{\frac{1}{12^2}+\frac{1}{16^2}+\frac{1}{\alpha^2}}}\]

 
 \[\sqrt{\frac{1}{12^2}+\frac{1}{16^2}+\frac{1}{\alpha^2}}=\frac{5}{4\alpha}\]

 
 \[\frac{12^2+16^2}{12^2\cdot16^2}+\frac{1}{\alpha^2}=\frac{25}{16\alpha^2}\]

 
 \[\frac{20^2}{192^2}=\frac{9}{16\alpha^2}\]

 
 \[\frac{20}{192}=\frac{3}{4\alpha}\]

 
 \[\alpha=\frac{3}{4}\cdot\frac{192}{20}=\frac{36}{5}\]

 Thus the answer is $36+5=\boxed{041}$.

 





\textbf{Solution 8 (Similar to Solution 5)}

Again, consider the total height as $2x$, which leads to the side lengths of the triangle being $10$, $\sqrt{x^2 + 64}$, and $\sqrt{x^2 + 36}$. Consider the angle $\theta$ between the two sides of $\sqrt{x^2+36}$ and $10$. 

 Since the area of a triangle is $\frac{1}{2}ab\sin c$, we can find that $\frac{1}{2}(10)\sqrt{x^2+36}\sin \theta = 30$. By cosine law, $x^2+64 = x^2+36+100-2(10)\sqrt{x^2+36}\cos \theta$, which simplifies to $20\sqrt{x^2+36}\cos \theta = 72$.

 Divide the first equation by the second to get $\frac{\sin\theta}{\cos\theta} = \frac{5}{3}$, and solve to find $\sin\theta=\frac{5}{\sqrt{34}}$. Substitute this back into the first equation to get the following:
 \[\sqrt{x^2+36} = \frac{6}{5}\sqrt{34}\]
 \[x^2+36 = \frac{36}{25}(34)\]
 \[x^2 = \frac{36(34)-36(25)}{25}\]
 \[x = \frac{18}{5}\]
 \[2x = \frac{36}{5}\].

 Thus, our answer is $\boxed{041}$.

 




	\item (\textit{AIME I 2013}) The domain of the function $f(x) = \arcsin(\log_{m}(nx))$ is a closed interval of length $\frac{1}{2013}$ , where $m$ and $n$ are positive integers and $m>1$. Find the remainder when the smallest possible sum $m+n$ is divided by $1000$.


 
 \textbf{Answer: }371



\textbf{Solution 1}


 \[x \in \text{Dom}(f)\] 
 \[\iff \log_{m}(nx) \in [-1, 1]\] 
 \[\iff nx \in [\frac{1}{m}, m]\] 
 \[\iff x \in [\frac{1}{mn}, \frac{m}{n}]\] 
 \[\implies \frac{1}{2013} = \frac{m}{n}-\frac{1}{mn} \implies n = \frac{2013(m^2-1)}{m}\]The last function of $m$ is strictly increasing over $[1, \infty)$, so we want the minimal $m$. 

 $m=2$ doesn't work, giving a non-integer value of $n$. But $m=3$ does, yielding $n = 5368$, so $m+n =5\boxed{371}$

 -sillybone

 





\textbf{Solution 2}

We know that the domain of $\text{arcsin}$ is $[-1, 1]$, so $-1 \le \log_m nx \le 1$. Now we can apply the definition of logarithms:
 \[m^{-1} = \frac1m \le nx \le m\] 
 \[\implies \frac{1}{mn} \le x \le \frac{m}{n}\]Since the domain of $f(x)$ has length $\frac{1}{2013}$, we have that
 \[\frac{m}{n} - \frac{1}{mn} = \frac{1}{2013}\] 
 \[\implies \frac{m^2 - 1}{mn} = \frac{1}{2013}\]

 A larger value of $m$ will also result in a larger value of $n$ since $\frac{m^2 - 1}{mn} \approx  \frac{m^2}{mn}=\frac{m}{n}$ meaning $m$ and $n$ increase about linearly for large $m$ and $n$. So we want to find the smallest value of $m$ that also results in an integer value of $n$. The problem states that $m > 1$. Thus, first we try $m = 2$:
 \[\frac{3}{2n} = \frac{1}{2013} \implies 2n = 3 \cdot 2013 \implies n \notin \mathbb{Z}\]  Now, we try $m=3$:
 \[\frac{8}{3n} = \frac{1}{2013} \implies 3n = 8 \cdot 2013 \implies n = 8 \cdot 671 = 5368\]Since $m=3$ is the smallest value of $m$ that results in an integral $n$ value, we have minimized $m+n$, which is $5368 + 3 = 5371 \equiv \boxed{371} \pmod{1000}$.

 





\textbf{Solution 3}

We start with the same method as above. The domain of the arcsin function is $[-1, 1]$, so $-1 \le \log_{m}(nx) \le 1$. 

 
 \[\frac{1}{m} \le nx \le m\] 
 \[\frac{1}{mn} \le x \le \frac{m}{n}\] 
 \[\frac{m}{n} - \frac{1}{mn} = \frac{1}{2013}\] 
 \[n = 2013m - \frac{2013}{m}\]

 For $n$ to be an integer, $m$ must divide $2013$, and $m > 1$. To minimize $n$, $m$ should be as small as possible because increasing $m$ will decrease $\frac{2013}{m}$, the amount you are subtracting, and increase $2013m$, the amount you are adding; this also leads to a small $n$ which clearly minimizes $m+n$.

 We let $m$ equal $3$, the smallest factor of $2013$ that isn't $1$. Then we have $n = 2013*3 - \frac{2013}{3} = 6039 - 671 = 5368$

 $m + n = 5371$, so the answer is $\boxed{371}$.

 





\textbf{Solution 4 (Operation Quadratics)}

Note that we need $-1\le f(x)\le 1$, and this eventually gets to $\frac{m^2-1}{mn}=\frac{1}{2013}$. From there, break out the quadratic formula and note that 
 \[m= \frac{n+\sqrt{n^2+4026^2}}{2013\times 2}.\] Then we realize that the square root, call it $a$, must be an integer. Then $(a-n)(a+n)=4026^2.$ 

 Observe carefully that $4026^2 = 2\times 2\times 3\times 3\times 11\times 11\times 61\times 61$! It is not difficult to see that to minimize the sum, we want to minimize $n$ as much as possible. Seeing that $2a$ is even, we note that a $2$ belongs in each factor. Now, since we want to minimize $a$ to minimize $n$, we want to distribute the factors so that their ratio is as small as possible (sum is thus minimum). The smallest allocation of $2, 61, 61$ and $2, 11, 3, 3, 11$ fails; the next best is $2, 61, 11, 3, 3$ and $2, 61, 11$, in which $a=6710$ and $n=5368$. That is our best solution, upon which we see that $m=3$, thus $\boxed{371}$.

 





\textbf{Solution 5 (Factor! and slightly cheese)}

By considering the range of $\arcsin$, we obtain $-1 \leq\log_{m} (nx) \leq 1 \Rightarrow \frac{1}{mn} \leq x \leq \frac{m}{n}$. Hence we consider the difference: $\frac{m}{n} - \frac{1}{mn} = \frac{m^2 - 1}{mn} = \frac{1}{2013}$. Now just cross multiply to obtain $2013(m-1)(m+1) = mn$, and that $2013 = 3 \cdot 11 \cdot 61$. We could try a few cases: $m = 3$, for instance, yields $n = 671 \cdot 8 = 5368$. Adding $3$ gives $5371$. Now, if we consider greater $m$, we see that the term $(m-1)(m+1)$ grows very fast: this is also due to the fact that $61,11 \gg 3$. Hence, we may just conclude that the answer is $5371 \Rightarrow \boxed{371}$.

 




	\item (\textit{AIME I 2013}) A paper equilateral triangle $ABC$ has side length $12$. The paper triangle is folded so that vertex $A$ touches a point on side $\overline{BC}$ a distance $9$ from point $B$. The length of the line segment along which the triangle is folded can be written as $\frac{m\sqrt{p}}{n}$, where $m$, $n$, and $p$ are positive integers, $m$ and $n$ are relatively prime, and $p$ is not divisible by the square of any prime. Find $m+n+p$.


 
                    \begin{figure}[H]
                        \centering
                        \includegraphics[width=0.5\textwidth]{Problem/AIME_2013_I_9_Problem_Figure_1.png}
                    \end{figure}
                    


 
 \textbf{Answer: }113



\textbf{Solution 1}

Let $M$ and $N$ be the points on $\overline{AB}$ and $\overline{AC}$, respectively, where the paper is folded. Let $D$ be the point on $\overline{BC}$ where the folded $A$ touches it.
                    \begin{figure}[H]
                        \centering
                        \includegraphics[width=0.5\textwidth]{Solution/AIME_2013_I_9_Solution_Figure_1.png}
                    \end{figure}
                    We have $AF=6\sqrt{3}$ and $FD=3$, so $AD=3\sqrt{13}$. Denote $\angle DAF = \theta$; we get $\cos\theta = 2\sqrt{3}/\sqrt{13}$. 

 In triangle $AXY$, $AY=\frac 12 AD = \frac 32 \sqrt{13}$, and $AX=AY\sec\theta =\frac{13}{4}\sqrt{3}$. 

 In triangle $AMX$, we get $\angle AMX=60^\circ-\theta$ and then use sine-law to get $MX=\frac 12 AX\csc(60^\circ-\theta)$; similarly, from triangle $ANX$ we get $NX=\frac 12 AX\csc(60^\circ+\theta)$. Thus 
 \[MN=\frac 12 AX(\csc(60^\circ-\theta) +\csc(60^\circ+\theta)).\] Since $\sin(60^\circ\pm \theta) = \frac 12 (\sqrt{3}\cos\theta \pm \sin\theta)$, we get
 \begin{align*}     \csc(60^\circ-\theta) +\csc(60^\circ+\theta) &= \frac{\sqrt{3}\cos\theta}{\cos^2\theta - \frac 14} = \frac{24 \cdot \sqrt{13}}{35} \end{align*}
Then 
 \[MN = \frac 12 AX \cdot \frac{24 \cdot \sqrt{13}}{35} = \frac{39\sqrt{39}}{35}\]

 The answer is $39 + 39 + 35 = \boxed{113}$.

 





\textbf{Solution 2}

Let $P$ and $Q$ be the points on $\overline{AB}$ and $\overline{AC}$, respectively, where the paper is folded.

 Let $D$ be the point on $\overline{BC}$ where the folded $A$ touches it.

 Let $a$, $b$, and $x$ be the lengths $AP$, $AQ$, and $PQ$, respectively.

 We have $PD = a$, $QD = b$, $BP = 12 - a$, $CQ = 12 - b$, $BD = 9$, and $CD = 3$.

 Using the Law of Cosines on $BPD$:

 $a^{2} = (12 - a)^{2} + 9^{2} - 2 \times (12 - a) \times 9 \times \cos{60}$

 $a^{2} = 144 - 24a + a^{2} + 81 - 108 + 9a$

 $a = \frac{39}{5}$

 Using the Law of Cosines on $CQD$:

 $b^{2} = (12 - b)^{2} +3^{2} - 2 \times (12 - b) \times 3 \times \cos{60}$

 $b^{2} = 144 - 24b + b^{2} + 9 - 36 + 3b$

 $b = \frac{39}{7}$

 Using the Law of Cosines on $DPQ$:

 $x^{2} = a^{2} + b^{2} - 2ab \cos{60}$

 $x^{2} = (\frac{39}{5})^2 + (\frac{39}{7})^2 - (\frac{39}{5} \times \frac{39}{7})$

 $x = \frac{39 \sqrt{39}}{35}$

 The solution is $39 + 39 + 35 = \boxed{113}$.

 





\textbf{Solution 3}

Proceed with the same labeling as in Solution 1. 

 $\angle B = \angle C = \angle A = \angle PDQ = 60^\circ$

 $\angle PDB + \angle PDQ + \angle QDC = \angle QDC + \angle CQD + \angle C = 180^\circ$

 Therefore, $\angle PDB = \angle CQD$.

 Similarly, $\angle BPD = \angle QDC$.

 Now, $\bigtriangleup BPD$ and $\bigtriangleup CDQ$ are similar triangles, so

 $\frac{3}{12-a} = \frac{12-b}{9} = \frac{b}{a}$.

 Solving this system of equations yields $a = \frac{39}{5}$ and $b = \frac{39}{7}$.

 Using the Law of Cosines on $APQ$:

 $x^{2} = a^{2} + b^{2} - 2ab \cos{60}$

 $x^{2} = (\frac{39}{5})^2 + (\frac{39}{7})^2 - (\frac{39}{5} \times \frac{39}{7})$

 $x = \frac{39 \sqrt{39}}{35}$

 The solution is $39 + 39 + 35 = \boxed{113}$.

 \textbf{Note} Once you find $DP$ and $DQ$, you can scale down the triangle by a factor of $\frac{39}{35}$ so that all sides are integers. Applying Law of cosines becomes easier, you just need to remember to scale back up.

 





\textbf{Solution 4 (Coordinate Bash)}

We let the original position of $A$ be $A$, and the position of $A$ after folding be $D$. Also, we put the triangle on the coordinate plane such that $A=(0,0)$, $B=(-6,-6\sqrt3)$, $C=(6,-6\sqrt3)$, and $D=(3,-6\sqrt3)$.

 
                    \begin{figure}[H]
                        \centering
                        \includegraphics[width=0.5\textwidth]{Solution/AIME_2013_I_9_Solution_Figure_1.png}
                    \end{figure}
                    

 Note that since $A$ is reflected over the fold line to $D$, the fold line is the perpendicular bisector of $AD$. We know $A=(0,0)$ and $D=(3,-6\sqrt3)$. The midpoint of $AD$ (which is a point on the fold line) is $(\frac32, -3\sqrt3)$. Also, the slope of $AD$ is $\frac{-6\sqrt3}{3}=-2\sqrt3$, so the slope of the fold line (which is perpendicular), is the negative of the reciprocal of the slope of $AD$, or $\frac{1}{2\sqrt3}=\frac{\sqrt3}{6}$. Then, using point slope form, the equation of the fold line is
 \[y+3\sqrt3=\frac{\sqrt3}{6}\left(x-\frac32\right)\]
 \[y=\frac{\sqrt3}{6}x-\frac{13\sqrt3}{4}\]Note that the equations of lines $AB$ and $AC$ are $y=\sqrt3x$ and $y=-\sqrt3x$, respectively. We will first find the intersection of $AB$ and the fold line by substituting for $y$:
 \[\sqrt3 x=\frac{\sqrt3}{6}x-\frac{13\sqrt3}{4}\]
 \[\frac{5\sqrt3}{6}x=-\frac{13\sqrt3}{4} \implies x=-\frac{39}{10}\]Therefore, the point of intersection is $\left(-\frac{39}{10},-\frac{39\sqrt3}{10}\right)$. Now, lets find the intersection with $AC$. Substituting for $y$ yields
 \[-\sqrt3 x=\frac{\sqrt3}{6}x-\frac{13\sqrt3}{4}\]
 \[\frac{-7\sqrt3}{6}x=-\frac{13\sqrt3}{4} \implies x=\frac{39}{14}\]Therefore, the point of intersection is $\left(\frac{39}{14},-\frac{39\sqrt3}{14}\right)$. Now, we just need to use the distance formula to find the distance between $\left(-\frac{39}{10},-\frac{39\sqrt3}{10}\right)$ and $\left(\frac{39}{14},-\frac{39\sqrt3}{14}\right)$.
 \[\sqrt{\left(\frac{39}{14}+\frac{39}{10}\right)^2+\left(-\frac{39\sqrt3}{14}+\frac{39\sqrt3}{10}\right)^2}\]The number 39 is in all of the terms, so let's factor it out:
 \[39\sqrt{\left(\frac{1}{14}+\frac{1}{10}\right)^2+\left(-\frac{\sqrt3}{14}+\frac{\sqrt3}{10}\right)^2}=39\sqrt{\left(\frac{6}{35}\right)^2+\left(\frac{\sqrt3}{35}\right)^2}\]
 \[\frac{39}{35}\sqrt{6^2+\sqrt3^2}=\frac{39\sqrt{39}}{35}\]Therefore, our answer is $39+39+35=\boxed{113}$, and we are done.

 Solution by nosaj.

 





\textbf{Solution 5}

Note: this requires lots of calculations that increase your chance of errors, but it only requires simple understanding of areas, similar triangles, and Heron's formula. I'll just put the strategy here because I am too lazy to calculate it myself right now. 

 1. Notice that the two triangles on the sides of the folded corner are similar. using this, we can find that the side lengths of them are $9,7.8,4.2$ and $3, \frac{45}{7}, \frac{39}{7}$

 2. Use heron's formula to find the areas of those two triangles. Remember that it is $\sqrt{s(s-a)(s-b)(s-c)}$

 3. Using the area of these triangles, we can find the area of the triangle with the length we need. 

 4. Use heron's formula again, with the unknown length as $x$, and since we know the area and the other two side lengths, we can just solve for x with this equation. 

 -EmilyQ, minor latex edits by Voidling

 





\textbf{Solution 6(Easy)}

Thanks to Solution 1 for the diagram below:

 
                    \begin{figure}[H]
                        \centering
                        \includegraphics[width=0.5\textwidth]{Solution/AIME_2013_I_9_Solution_Figure_1.png}
                    \end{figure}
                    

 We will use the notation already on the diagram, but our solution is slightly different. 

 We will only need M and N. 

 Let NC be length a, which implies NA be 12-a. 

 Also, AC = 3 because AB = 9

 By the Law of Cosines on NCA, 

 $(12-a)^2=a^2+3^2-2(a)(3)(cos60)$

 which simplifies to: 

 $a=\frac{45}{7}$

 Which means that NC = 45/7 and NA = 39/7. 

 We can do the same thing for MBA. 

 This time, MB = b. 

 $(12-b)^2=b^2+81-2(9)(b)(cos60)$

 Which gives:

 $b=\frac{21}{5}$Which implies that MA = $\frac{39}{5}$Now, since MAN is 60 degrees, we can apply the Law of Cosines again(I know, I don't like bashy things too) to get:

 $c^2=(\frac{39}{5})^2+39/7^2-2(39/7)(39/5)(cos60)$

 Which leads us to our answer => 113

 





\textbf{Solution 7 (LoC)}


                    \begin{figure}[H]
                        \centering
                        \includegraphics[width=0.5\textwidth]{Solution/AIME_2013_I_9_Solution_Figure_1.png}
                    \end{figure}
                    As shown in the diagram, let $MA'=x$. Since the triangle is equilateral, we have $BM=12-x$. Similarly, let $NA'=y$, so $NC=12-y$.

 Because $\angle MBA=\angle NCA=60^\circ$, we may apply the Law of Cosines in the form $a^2+b^2-ab=c^2$. (Note that we could also find these side lengths through similar triangles as mentioned in the previous solutions)

 Solving for $x$, we have
 \[(12-x)^2+9^2-9(12-x)=x^2.\]Simplifying,
 \[x^2-24x+144-27+9x=x^2,\]so
 \[15x=117 \quad \Rightarrow \quad x=\frac{39}{5}.\]

 Solving for $y$ similarly,
 \[(12-y)^2+3^2-3(12-y)=y^2.\]Simplifying,
 \[y^2-24y+144+9-36+3y=y^2,\]which gives
 \[21y=117 \quad \Rightarrow \quad y=\frac{39}{7}.\]

 


 Now apply the Law of Cosines again to find $MN$:
 \[MN^2=\left(\frac{39}{5}\right)^2+\left(\frac{39}{7}\right)^2-\left(\frac{39}{5}\right)\left(\frac{39}{7}\right).\]Factoring,
 \[MN^2=39^2\left(\frac{1}{25}+\frac{1}{49}-\frac{1}{35}\right)=\frac{39^3}{35^2}.\]Thus,
 \[MN=\frac{39\sqrt{39}}{35}.\]

 Since the problem asks for $m+n+p$ when the length is written as $\frac{m\sqrt{p}}{n}$, we have
 \[m+n+p=39+39+35=\boxed{113}.\]

 This video is private for some reason ~get-rickrolled

 




	\item (\textit{AIME I 2013}) There are nonzero integers $a$, $b$, $r$, and $s$ such that the complex number $r+si$ is a zero of the polynomial $P(x)={x}^{3}-a{x}^{2}+bx-65$. For each possible combination of $a$ and $b$, let ${p}_{a,b}$ be the sum of the zeros of $P(x)$. Find the sum of the ${p}_{a,b}$'s for all possible combinations of $a$ and $b$.


 



 
 \textbf{Answer: }080



\textbf{Solution}

Since $r+si$ is a root, by the Complex Conjugate Root Theorem, $r-si$ must be the other imaginary root. Using $q$ to represent the real root, we have

 $(x-q)(x-r-si)(x-r+si) = x^3 -ax^2 + bx -65$

 Applying difference of squares, and regrouping, we have

 $(x-q)(x^2 - 2rx + (r^2 + s^2)) = x^3 -ax^2 + bx -65$

 So matching coefficients, we obtain

 $q(r^2 + s^2) = 65$

 $b = r^2 + s^2 + 2rq$

 $a = q + 2r$

 By Vieta's each ${p}_{a,b} = a$ so we just need to find the values of $a$ in each pair.We proceed by determining possible values for $q$, $r$, and $s$ and using these to determine $a$ and $b$.

 If $q = 1$, $r^2 + s^2 = 65$so (r, s) = $(\pm1, \pm 8), (\pm8, \pm 1), (\pm4, \pm 7), (\pm7, \pm 4)$

 Similarly, for $q = 5$, $r^2 + s^2 = 13$ so the pairs $(r,s)$ are$(\pm2, \pm 3), (\pm3, \pm 2)$

 For $q = 13$, $r^2 + s^2 = 5$ so the pairs $(r,s)$ are$(\pm2, \pm 1), (\pm1, \pm 2)$

 Now we can disregard the plus minus signs for s because those cases are included as complex conjugates of the counted cases.The positive and negative values of r will cancel, so the sum of the  ${p}_{a,b} = a$ for $q = 1$ is $q$ times the number of distinct $r$ values (as each value of $r$ generates a pair $(a,b)$).Our answer is then $(1)(8) + (5)(4) + (13)(4) = \boxed{080}$.

 \textbf{Note} We consider all values of $q$ times their distinct $r$ values as $a = q+2r$ in which all $2r$ cancels out, so we just count the number of $q$. The wording in the solution was slightly jumpy.

 \textbf{Remark:} The complex conjugate theorem states that a polynomial with real coefficients must have an even amount of complex roots.One of them is the complex number $a + bi$, and the other is it's conjugate, or $a - bi$. These, when multiplied cancel out and become real numbers ($a^2 + b^2$). Similar logic for addition. 

 




	\item (\textit{AIME I 2013}) Ms. Math's kindergarten class has $16$ registered students. The classroom has a very large number, $N$, of play blocks which satisfies the conditions:


 (a) If $16$, $15$, or $14$ students are present in the class, then in each case all the blocks can be distributed in equal numbers to each student, and


 (b) There are three integers $0 < x < y < z < 14$ such that when $x$, $y$, or $z$ students are present and the blocks are distributed in equal numbers to each student, there are exactly three blocks left over.


 Find the sum of the distinct prime divisors of the least possible value of $N$ satisfying the above conditions.


 
 \textbf{Answer: }148



\textbf{Solution 1}

$N$ must be some multiple of $\text{lcm}(14, 15, 16)= 2^{4}\cdot 3\cdot 5\cdot 7$ ; this $lcm$ is hereby denoted $k$ and $N = qk$.

 $1$, $2$, $3$, $4$, $5$, $6$, $7$, $8$, $10$, and $12$ all divide $k$, so $x, y, z = 9, 11, 13$

 We have the following three modulo equations:

 $nk\equiv 3\pmod{9}$

 $nk\equiv 3\pmod{11}$

 $nk\equiv 3\pmod{13}$

 To solve the equations, you can notice the answer must be of the form $9\cdot 11\cdot 13\cdot m + 3$ where $m$ is an integer. 

 This must be divisible by $lcm$ $(14, 15, 16)$, which is $560\cdot 3$. 

 Therefore, $\frac{9\cdot 11\cdot 13\cdot m + 3}{560\cdot 3} = q$, which is an integer. Factor out $3$ and divide to get $\frac{429m+1}{560} = q$. Therefore, $429m+1=560q$. We can use \href{https://artofproblemsolving.com/wiki/index.php?title=Bezout\%27s\_Lemma}{Bezout's Identity} or a \href{https://artofproblemsolving.com/wiki/index.php?title=Euclidean\_algorithm}{Euclidean algorithm} bash to solve for the least of $m$ and $q$. 

 We find that the least $m$ is $171$ and the least $q$ is $131$. 

 Since we want to factor $1680\cdot q$, don't multiply: we already know that the prime factors of $1680$ are $2$, $3$, $5$, and $7$, and since $131$ is prime, we have $2 + 3 + 5 + 7 + 131 = \boxed{148}$.

 





\textbf{Solution 2}

Note that the number of play blocks is a multiple of the LCM of $16$, $15$, and $14$. The value of this can be found to be $(16)(15)(7) = 1680$. This number is also divisible by $1$, $2$, $3$, $4$, $5$, $6$, $7$, $8$, $10$, and $12$, thus, the three numbers $x, y, z$ are $9, 11, 13$.

 Thus, $1680k\equiv 3$ when taken mod $9$, $11$, $13$. Since $1680$ is congruent to $6$ mod $9$ and $3$ mod $13$, and congruent to $8$ mod $11$, the number $k$ must be a number that is congruent to $1$ mod $13$, $2$ mod $3$ (because $6$ is a multiple of $3$, which is a factor of $9$ that can be divided out) and cause $8$ to become $3$ when multiplied under modulo $11$.

 Looking at the last condition shows that $k\equiv 10$ mod $11$ (after a bit of bashing) and is congruent to $1$ mod $13$ and $2$ mod $3$ as previously noted. Listing out the numbers congruent to $10$ mod $11$ and $1$ mod $13$ yield the following lists:

 $10$ mod $11$: $21$, $32$, $43$, $54$, $65$, $76$, $87$, $98$, $109$, $120$, $131$...

 $1$ mod $13$: $14$, $27$, $40$, $53$, $66$, $79$, $92$, $105$, $118$, $131$, $144$, $157$, $170$...

 Both lists contain $x$ elements where $x$ is the modulo being taken, thus, there must be a solution in these lists as adding $11(13)$ to this solution yields the next smallest solution. In this case, $131$ is the solution for $k$ and thus the answer is $1680(131)$. Since $131$ is prime, the sum of the prime factors is $2 + 3 + 5 + 7 + 131 = \boxed{148}$.

 





\textbf{Solution 3}

It is obvious that $N=a\cdot 2^4 \cdot 3\cdot 5\cdot 7$ and so the only mod $3$ number of students are $9, 11, 13$. Therefore, $N=1287\cdot k+3$. Try some approaches and you will see that this one is one of the few successful ones:

 Start by setting the two $N$ equations together, then we get $1680a=1287k+3$. Divide by $3$. Note that since the RHS is $1\pmod{3}$, and since $560$ is $2\pmod{3}$, then $a=3b+2$, where $b$ is some nonnegative integer, because $a$ must be $2\pmod{3}$.

 This reduces to $560 \cdot 3b + 1119 = 429k$. Now, take out the $11!$ With the same procedure, $b=11c-1$, where $c$ is some nonnegative integer.

 You also get $c=13d+4$, at which point $k=171+560d$. $d$ cannot be equal to $0$. Therefore, $c=4, b=43, a=131$, and we know the prime factors of $N$ are $2, 3, 5, 7, 131$ so the answer is $\boxed{148}$.

 





\textbf{Solution 4}

We start by noticing that $N = a\textbf{lcm}(14, 15, 16) = 1680a$ for some integer $a$ in order to satisfy the first condition. 

 Next, we satisfy the second condition. Since $x<y<z < 14$ must leave a remainder when dividing $1680a$, they are not divisors of $1680x$. Thus, we can eliminate all $y \le 14$ s.t. $\gcd(y, 1680x) \ne 1$ which leaves $(x, y, z) = (9, 11, 13)$. Thus, $N = 1680a \equiv 3 \pmod 9 \equiv 3 \pmod {11} \equiv 3 \pmod {13}$. Now, we seek to find the least $a$ which satisfies this set of congruences.

 By Chinese Remainder Theorem on the first two congruences, we find that $a \equiv 32 \pmod {33}$ (we divide by three before proceeding in the first congruence to ensure the minimal solution). Finally, by CRT again on $a \equiv 32 \pmod {33}$ and $1680a \equiv 3 \pmod {13}$ we find that $a \equiv 131 \pmod {429}$. 

 Thus, the minimal value of $N$ is possible at $a = 131$. The prime factorization of this minimum value is $2^4 \cdot 3 \cdot 5 \cdot 7 \cdot 131$ and so the answer is $2 + 3 + 5 + 7 + 131 = \boxed{148}$.

 


 





\textbf{Solution 5}

As the problem stated, the number of boxes is definitely a multiple of $lcm(14,15,16)=1680$, so we assume total number of boxes is $1680k$

 Then, according to $(b)$ statement, we get $1680k \equiv 3 \pmod x \equiv 3 \pmod {y} \equiv 3 \pmod {z}$. So we have $lcm(x,y,z)+3+m\cdot lcm(x,y,z)=1680k$, we just write it to be $(1+n)lcm(x,y,z)=1680k-3$ Which tells that $x,y,z$ must be all odd number. Moreover, we can see $(1+m)lcm(x,y,z)$ can't be a multiple of $3,5,7$(as $1680$ is a multiple of $5,7$) which means that $lcm(x,y,z)=lcm(9,11,13)=1287$ We let $n=1+m$

 Now, we write $1287n+3=1680k, 429n+1=560k$It is true that $n\equiv 1 \pmod{10}$, let $n=10p+1$, it has $429p+43=56k$ Then, $p$ must be odd, let $p=2q+1$, it indicates $429q+236=28k$ Now, $q$ must be even, $q=2s$ tells $429s+118=14k$ Eventually, $s$ must be even, $s=2y$, $858y+59=7k$, $y=1$ is the smallest. This time, $k=131$

 So the number of balls is $1680\cdot 131=2^4\cdot 3\cdot 5\cdot 7 \cdot 131$, the desired value is $2+3+5+7+131=\boxed{148}$

 


 





\textbf{Solution 6(CRT Bash)}

From part (a), we know that $2^4\cdot3\cdot5\cdot7 | N$. From part (b), we know that $N \equiv 3 \pmod {1287}$. We can expand on part (a) by saying that $N = 1680k$ for some $k$. Rather than taking the three modulos together, we take them individually. 
 \[1680k \equiv 6k \equiv 3 \pmod 9\]
 \[k \equiv 2^{-1} \pmod 9\]The inverse of 2 mod 9 is easily seen to be $5$. 
 \[k \equiv 5 \pmod 9\]

 Now moving to the second modulo which we leave as follows,
 \[1680k \equiv 8k \equiv 3 \pmod {11}\]Now the last modulo,
 \[1680k \equiv 3k \equiv 3 \pmod {13}\]
 \[k \equiv 1 \pmod{13}\]CRT on the first and the third one results in $k \equiv 4 \pmod {117}$. Now doing the second one and the one we just made, $k \equiv 131 \pmod{1287}$. Thus, the smallest value that works for $k = 131$. Thus $N = 2^4\cdot3\cdot5\cdot7\cdot131$ $2+3+5+7+131 = \boxed{148}$

 





\textbf{Solution 7 (similar to solution 6)}

Define $N_b$ to be the base value of $N$. In other terms, we call $N_b$ the "skeleton". We proceed with an archeological approach.

 Firstly, $N_b \equiv 0 \pmod{16} \equiv 0 \pmod{15} \equiv 0 \pmod{14}$ which implies $N_b \equiv 0 \pmod{2^4 \cdot 3 \cdot 5 \cdot 7} \equiv N_b \equiv 0 \pmod{1680}$.

 Notice this is valid for all integers besides $p \in \{9,11,13 \}$. We require $N-3$ to be divisible by ALL integers less than or equal to 16. 

 So, $N_b \ | \ N$ and $N \equiv 3 \pmod{9}, N \equiv 3 \pmod{11}, N \equiv 3 \pmod{13}, N \equiv 0 \pmod{1680}$. Then, $N \equiv 3 \pmod{1287}$. Then, $1287k \equiv 1677 \pmod{1680}$. Dividing everything by 3, we have $429k \equiv 559 \pmod{560}$. We solve via modular inverse rather than CRT, and the modular inverse of 429 through Bézout’s (or extended Euclidean) is 389. So, $k \equiv 171 \pmod{560}$. Substitution, we have $N = 3 + 1287(171)$ or $N = 220080$. The prime factorization of this is $N = 2^4\cdot3\cdot5\cdot7\cdot131$, and the value is $2+3+5+7+131= \boxed{148}$.

 




	\item (\textit{AIME I 2013}) Let $\bigtriangleup PQR$ be a triangle with $\angle P = 75^o$ and $\angle Q = 60^o$. A regular hexagon $ABCDEF$ with side length 1 is drawn inside $\triangle PQR$ so that side $\overline{AB}$ lies on $\overline{PQ}$, side $\overline{CD}$ lies on $\overline{QR}$, and one of the remaining vertices lies on $\overline{RP}$. There are positive integers $a, b, c,$ and $d$ such that the area of $\triangle PQR$ can be expressed in the form $\frac{a+b\sqrt{c}}{d}$, where $a$ and $d$ are relatively prime, and c is not divisible by the square of any prime. Find $a+b+c+d$.


 
 \textbf{Answer: }021



\textbf{Solution 1}

First, find that $\angle R = 45^\circ$.Draw $ABCDEF$. Now draw $\bigtriangleup PQR$ around $ABCDEF$ such that $Q$ is adjacent to $C$ and $D$. The height of $ABCDEF$ is $\sqrt{3}$, so the length of base $QR$ is $2+\sqrt{3}$. Let the equation of $\overline{RP}$ be $y = x$. Then, the equation of $\overline{PQ}$ is $y = -\sqrt{3} (x - (2+\sqrt{3})) \to y = -x\sqrt{3} + 2\sqrt{3} + 3$. Solving the two equations gives $y = x = \frac{\sqrt{3} + 3}{2}$. The area of $\bigtriangleup PQR$ is $\frac{1}{2} * (2 + \sqrt{3}) * \frac{\sqrt{3} + 3}{2} = \frac{5\sqrt{3} + 9}{4}$. $a + b + c + d = 9 + 5 + 3 + 4 = \boxed{021}$

 


 Note (different ending): When you have the length of the base $QR$, you don't need to find the equations of the lines $QP$ and $PR$. Instead, make an altitude from $P$ to $QR$, and call the foot $M$. $QPM$ is a $45, 45, 90$ triangle and $PMR$ is a $30, 60, 90$ triangle. And, they both share $PM$. So, we can set $RM$ as $x$, so $PM$ is $x\sqrt{3}$. Since $QPM$ is a $45, 45, 90$ triangle, $PM=MQ=x\sqrt{3}$. The base $QR$ can be written as $QM+MR=x+x\sqrt{3}=2+\sqrt{3}$. Solve this equation and $x=\frac{\sqrt{3}+1}{2}$ and $PM=x\sqrt{3}=\frac{\sqrt{3}+3}{2}$. Multiply this by base $QR$ and divide by $2$ to get the area of triangle $PQR$ which is $\frac{9+5\sqrt{3}}{4}$. So, $a + b + c + d = 9 + 5 + 3 + 4 = \boxed{021}$-hwan

 





\textbf{Solution 2 (Cartesian Variation)}

Use coordinates. Call $Q$ the origin and $QP$ be on the x-axis. It is easy to see that $F$ is the vertex on $RP$. After labeling coordinates (noting additionally that $QBC$ is an equilateral triangle), we see that the area is $QP$ times $0.5$ times the coordinate of $R$. Draw a perpendicular of $F$, call it $H$, and note that $QP = 1 + \sqrt{3}$ after using the trig functions for $75$ degrees.

 Now, get the lines for $QR$ and $RP$: $y=\sqrt{3}x$ and $y=-(2+\sqrt{3})x + (5+\sqrt{3})$, whereupon we get the ordinate of $R$ to be $\frac{3+2\sqrt{3}}{2}$, and the area is $\frac{5\sqrt{3} + 9}{4}$, so our answer is $\boxed{021}$.

 





\textbf{Solution 3 (Trig)}

Angle chasing yields that both triangles $PAF$ and $PQR$ are $75$-$60$-$45$ triangles. First look at triangle $PAF$. Using Law of Sines, we find:

 $\frac{\frac{\sqrt{6} + \sqrt{2}}{4}}{1} = \frac{\frac{\sqrt{2}}{2}}{PA}$

 Simplifying, we find $PA = \sqrt{3} - 1$.Since $\angle{Q} = 60^\circ$, WLOG assume triangle $BQC$ is equilateral, so $BQ = 1$. So $PQ = \sqrt{3} + 1$.

 Apply Law of Sines again, 

 $\frac{\frac{\sqrt{2}}{2}}{\sqrt{3} + 1} = \frac{\frac{\sqrt{3}}{2}}{PR}$

 Simplifying, we find $PR = \frac{\sqrt{6}}{2} \cdot (1 + \sqrt{3})$. 

 $[PQR] = \frac{1}{2} \cdot PQ \cdot PR \cdot \sin 75^\circ$. 

 Evaluating and reducing, we get $\frac{9 + 5\sqrt{3}}{4},$thus the answer is $\boxed{021}$

 





\textbf{Solution 4 (Special Triangles)}


                    \begin{figure}[H]
                        \centering
                        \includegraphics[width=0.5\textwidth]{Solution/AIME_2013_I_12_Solution_Figure_1.png}
                    \end{figure}
                    

 As we can see, the $75^\circ$ angle of $\angle P$ can be split into a $45^\circ$ angle and a $30^\circ$ angle. This allows us to drop an altitude from point $P$ for $\triangle RPQ$ which intersects $\overline{AF}$ at point $a$ and $\overline{RQ}$ at point $b$. The main idea of our solution is to obtain enough sides of $\triangle RPQ$ that we are able to directly figure out its area (specifically by figuring out side $\overline{RQ}$ and $\overline{Pb}$). 

 We first begin by figuring out the length of $\overline{PQ}$. This can be easily done, since $\overline{AB}$ is simply $1$ (given in the problem) and $\overline{BQ}=1$ because $\triangle BCQ$ is an equilateral after some simple angle calculations. Now we need to find $\overline{PA}$. This is when we bring in some simple algebra. 

 PREPARATION:$\overline{aF}=\overline{Pa}$ (45-45-90 Right Triangle)

 $\overline{Pa}=\sqrt{3}\overline{Aa}$ (30-60-90 Right Triangle)

 $\overline{PA}=2\overline{Aa}$

 $\overline{Aa}+\overline{aF}=1$

 SOLVING:$\overline{Aa}+\sqrt{3}\overline{Aa}=1$

 so $\overline{Aa}(\sqrt{3}+1)=1$

 $\overline{Aa}=\frac{1}{\sqrt{3}+1}=\frac{\sqrt{3}-1}{2}$

 Finally, $\overline{PA}=2\cdot\frac{\sqrt{3}-1}{2}=\sqrt{3}-1$

 
Now, we can finally get the length of $\overline{PQ}$ by adding up $\overline{PA}+\overline{AB}+\overline{BQ}$, which is simply $(\sqrt{3}-1)+(1)+(1)=\sqrt{3}+1$

 
To get $\overline{RQ}$ and $\overline{Pb}$, we first work bit by bit.

 $\overline{Qb}=\frac{\overline{PQ}}{2}=\frac{\sqrt{3}+1}{2}$ (30-60-90 Right Triangle)

 $\overline{Pb}=\sqrt{3}\overline{Qb}=\frac{\sqrt{3}+3}{2}$ (same 30-60-90 Right Triangle)

 Since $\overline{Pb}=\overline{Rb}$ because of 45-45-90 right triangles, 

 $\overline{Rb}=\frac{\sqrt{3}+3}{2}$ too.

 Now, we can finally calculate $\overline{RQ}$, and it is $\overline{Rb}+\overline{Qb}=\frac{\sqrt{3}+3}{2}+\frac{\sqrt{3}+1}{2}=\sqrt{3}+2$.

 Finally, the area of $\triangle PRQ$ can be calculated by $\frac{1}{2}\cdot\overline{RQ}\cdot\overline{Pb}$, which is equal to $[\triangle PRQ]=\frac{1}{2} \cdot (\sqrt{3}+2) \cdot \frac{\sqrt{3}+3}{2} =\frac{9+5\sqrt{3}}{4}$. So the final answer is $9+5+3+4=\fbox{021}$.

 -by What do Humanitarians Eat?

 





\textbf{Solution 5 (Trig)}


                    \begin{figure}[H]
                        \centering
                        \includegraphics[width=0.5\textwidth]{Solution/AIME_2013_I_12_Solution_Figure_1.png}
                    \end{figure}
                    

 With some simple angle chasing we can show that $\triangle OJL$ and $\triangle MPL$ are congruent. This means we have a large equilateral triangle with side length $3$ and quadrilateral $OJQN$. We know that $[OJQN] = [\triangle NQL] - [\triangle OJL]$. Using Law of Sines and the fact that $\angle N = 45^{\circ}$ we know that $\overline{NL} = \sqrt{6}$ and the height to that side is $\frac{\sqrt{3} -1}{\sqrt{2}}$ so $[\triangle NQL] = \frac{3-\sqrt{3}}{2}$. Using an extremely similar process we can show that $\overline{OJ} = 2-\sqrt{3}$ which means the height to $\overline{LJ}$ is $\frac{2\sqrt{3}-3}{2}$. So the area of $\triangle OJL = \frac{2\sqrt{3}-3}{4}$. This means the area of quadrilateral $OJQN = \frac{3-\sqrt{3}}{2} - \frac{2\sqrt{3}-3}{4} = \frac{9-4\sqrt{3}}{4}$. So the area of our larger triangle is $\frac{9-4\sqrt{3}}{4} + \frac{9\sqrt{3}}{4} = \frac{9+5\sqrt{3}}{4}$. Therefore $9+5+3+4=\fbox{021}$.

 





\textbf{Solution 6 (Elementary Geo)}

We can find that $AF || CD || QR$. This means that the perpendicular from $P$ to $QR$ is perpendicular to $AF$ as well, so let that perpendicular intersect $AF$ at $G$, and the perpendicular intersect $QR$ at $H$. Set $AP=x$. Note that $\angle {PAG} = 60^\circ$, so $AG=\frac{x}{2}$ and $PG = GF = \frac{x\sqrt3}{2}$. Also, $1=AF=AG+GF=\frac{x}{2} + \frac{x\sqrt{3}}{2}$, so $x=\sqrt{3} - 1$. It's easy to calculate the area now, because the perpendicular from $P$ to $QR$ splits $\triangle{PQR}$ into a $30-60-90$ (PHQ) and a $45-45-90$ (PHR). From these triangles' ratios, it should follow that $QH=\frac{\sqrt{3} + 1}{2}, PH=HR=\frac{\sqrt{3}+3}{2}$, so the area is $\frac{1}{2} * PH * QR = \frac{1}{2} * PH * (QH + HR) = \frac{1}{2} * \frac{\sqrt{3} + 3}{2} * \frac{2\sqrt{3}+4}{2} = \boxed{\frac{9+5\sqrt{3}}{4}}$. $9+5+3+4=021$.By Mathscienceclass

 





\textbf{Solution 7 (Combination of 1 & 2)}

We can observe that $RD=DF$ (because $\angle R$ & $\angle RFD$ are both $45^\circ$). Thus we know that $RD$ is equivalent to the height of the hexagon, which is $\sqrt3$. Now we look at triangle $\triangle AFP$ and apply the Law of Sines to it. $\frac{1}{\sin{75}}=\frac{AP}{\sin{45}}$. From here we can solve for $AP$ and get that $AP=\sqrt{3}-1$. Now we use the Sine formula for the area of a triangle with sides $RQ$, $PQ$, and $\angle {RQP}$ to get the answer. Setting $PQ=\sqrt{3}+1$ and $QR=\sqrt{3}+2$ we get the expression $\frac{(\sqrt{3}+1)(\sqrt{3}+2)(\frac{\sqrt{3}}{2})}{2}$ which is  $\frac{9 + 5\sqrt{3}}{4}$. Thus our final answer is $9+5+3+4=\fbox{021}$.By AwesomeLife\_Math

 





\textbf{Solution 8 (Area Ratios + Trig)}

Note that $\triangle PAF \sim \triangle PQR$. Let $X$ be the foot of an altitude dropped from $F$ to $PQ$. Using trigonometry, we find that $PA = \sqrt3 - 1$ and $FX = \frac{\sqrt3}{2}$, thus the area of $\triangle PAF$ is $\frac{3 - \sqrt3}{4}$. Since $QA$ is clearly $2$, the scale factor between $\triangle PAF$ and $\triangle PQR$ is $\frac{\sqrt3 + 1}{\sqrt3 - 1}$, and thus $[PQR] = \frac{3 - \sqrt3}{4} \cdot \left( \frac{\sqrt3 + 1}{\sqrt3 - 1} \right)^2 = \frac{9 + 5\sqrt3}{4}$, and we extract $\boxed{021}$. ~ostriches88

 





\textbf{Solution 9 (I'm not sure if this has been posted before)}

Let $F$ be the point on $RP$. Drop the perpendicular $FD$, and notice the length of this is $\sqrt{3}$. Consider the $45-45-90$ triangle $\triangle FDR$. The area is clearly $\frac{3}{2}$. This also takes away a sixth of the original hexagon, so we add $\frac{5}{6} \cdot \frac{1^2 \sqrt{3}}{4}$. The equilateral triangle $\triangle BCQ$ has area $\frac{3}{4}$. Finally, the upper triangle, $\triangle PAF$, has $1$ side facing the $75^{\circ}$ angle, so consequently by law of sines we obtain the side opposing the $45$ degree angle has side length $\sqrt{3} - 1$ (recall that $\sin 75^{\circ} = \frac{\sqrt{6} + \sqrt{2}}{4}$). Then, we have $\frac{1}{2} (\sqrt{3} - 1)(1) \sin(60^{\circ}) = \frac{3-\sqrt{3}}{4}$. Adding the four areas together, we obtain $\frac{9 + 5\sqrt{3}}{4}$. (Don't forget that the area of the 45-45-90 triangle is not in fact $3$, unless you want a final answer of $27$ instead of the actual $\boxed{21}$).

 




	\item (\textit{AIME I 2013}) Triangle $AB_0C_0$ has side lengths $AB_0 = 12$, $B_0C_0 = 17$, and $C_0A = 25$. For each positive integer $n$, points $B_n$ and $C_n$ are located on $\overline{AB_{n-1}}$ and $\overline{AC_{n-1}}$, respectively, creating three similar triangles $\triangle AB_nC_n \sim \triangle B_{n-1}C_nC_{n-1} \sim \triangle AB_{n-1}C_{n-1}$. The area of the union of all triangles $B_{n-1}C_nB_n$ for $n\geq1$ can be expressed as $\frac pq$, where $p$ and $q$ are relatively prime positive integers. Find $q$.


 
 \textbf{Answer: }961



\textbf{Solution 1}

Note that every $B_nC_n$ is parallel to each other for any nonnegative $n$. Also, the area we seek is simply the ratio $k=\frac{[B_0B_1C_1]}{[B_0B_1C_1]+[C_1C_0B_0]}$, because it repeats in smaller and smaller units. Note that the area of the triangle, by Heron's formula, is 90.

 For ease, all ratios I will use to solve this problem are with respect to the area of $[AB_0C_0]$. For example, if I say some area has ratio $\frac{1}{2}$, that means its area is 45.

 Now note that $k=$ 1 minus ratio of $[B_1C_1A]$ minus ratio $[B_0C_0C_1]$. We see by similar triangles given that ratio $[B_0C_0C_1]$ is $\frac{17^2}{25^2}$. Ratio $[B_1C_1A]$ is $(\frac{336}{625})^2$, after seeing that $C_1C_0 = \frac{289}{625}$, . Now it suffices to find 90 times ratio $[B_0B_1C_1]$, which is given by 1 minus the two aforementioned ratios. Substituting these ratios to find $k$ and clearing out the $5^8$, we see that the answer is $90\cdot \frac{5^8-336^2-17^2\cdot 5^4}{5^8-336^2}$, which gives $q= \boxed{961}$.

 





\textbf{Solution 2}

Using Heron's Formula we can get the area of the triangle $\Delta AB_0C_0 = 90$. 

 Since $\Delta AB_0C_0 \sim \Delta B_0C_1C_0$ then the scale factor for the dimensions of  $\Delta B_0C_1C_0$ to $\Delta AB_0C_0$ is $\dfrac{17}{25}.$ 

 Therefore, the area of $\Delta B_0C_1C_0$ is $(\dfrac{17}{25})^2(90)$. Also, the dimensions of the other sides of the $\Delta B_0C_1C_0$ can beeasily computed: $\overline{B_0C_1}= \dfrac{17}{25}(12)$ and $\overline{C_1C_0} = \dfrac{17^2}{25}$. This allows us to compute one side of the triangle $\Delta AB_0C_0$, $\overline{AC_1} = 25 - \dfrac{17^2}{25} = \dfrac{25^2 - 17^2}{25}$. Therefore, the scale factor $\Delta AB_1C_1$ to $\Delta AB_0C_0$ is $\dfrac{25^2 - 17^2}{25^2}$ , which yields the length of $\overline{B_1C_1}$ as $\dfrac{25^2 - 17^2}{25^2}(17)$. Therefore, the scale factor for $\Delta B_1C_2C_1$ to $\Delta B_0C_1C_0$ is $\dfrac{25^2 - 17^2}{25^2}$. Some more algebraic manipulation will show that  $\Delta B_nC_{n+1}C_n$ to $\Delta B_{n-1}C_nC_{n-1}$ is still $\dfrac{25^2 - 17^2}{25^2}$.  Also, since the triangles are disjoint, the area of the union is the sum of the areas. Therefore, the area is the geometric series $\dfrac{90 \cdot 17^2}{25^2} \sum_{n=0}^{\infty} (\dfrac{25^2-17^2}{25^2})^2$At this point, it may be wise to "simplify" $25^2 - 17^2 = (25-17)(25+17) = (8)(42) = 336$.So the geometric series converges to $\dfrac{90 \cdot 17^2}{25^2} \dfrac{1}{1 - \dfrac{336^2}{625^2}} = \dfrac{90 \cdot 17^2}{25^2} \dfrac{625^2}{625^2 - 336^2}$.Using the difference of squares, we get $\dfrac{90 \cdot 17^2}{25^2}\dfrac{625^2}{(625 - 336)(625 + 336)}$, which simplifies to  $\dfrac{90 \cdot 17^2}{25^2} \dfrac{625^2}{(289)(961)}$. Cancelling all common factors, we get the reduced fraction  $= \dfrac{90 \cdot 25^2}{31^2}$. So $\frac{p}{q}=1-\frac{90 \cdot 25^2}{31^2}=\frac{90 \cdot 336}{961}$, yielding the answer $\fbox{961}$.

 





\textbf{Solution 3}

For this problem, the key is to find the $\frac{[\triangle{B_nB_{n-1}C_n}]}{[\triangle{AB_{n-1}C_{n-1}}]}$.

 The area of the biggest triangle is $90$ according to the Heron's formula easily

 Firstly, we discuss the ratio of  $\frac{[\triangle{B_0C_1C_0}]}{[\triangle{AB_0C_0}]}$ 

 Since the problem said that two triangles are similar, so $\frac{C_1C_0}{B_0C_0}=\frac{17}{25}$,

 Getting that $C_1C_0=\frac{289}{25}$, which is not hard to find that $AC_1=\frac{336}{25}$, Since $\frac{AB_1}{AB_0}=\frac{AC_1}{AC_0}=\frac{336}{625}$,

 we can find the ratio of $\frac{[\triangle{B_0B_1C_1}]}{[\triangle{AB_0C_0}]}=\frac{336}{625}\cdot\frac{289}{625}$, the common ratio between two similar triangles is $(\frac{336}{625})^2$, the similar triangles means two consecutive $(\triangle{AB_nC_n});(\triangle{AB_{n+1}C_{n+1}})$

 Now the whole summation of $S=1+(\frac{336}{625})^2+(\frac{336}{625})^4+....+(\frac{336}{625})^n=\frac{625^2}{961\cdot289}$

 The desired answer is $90\cdot\frac{336\cdot289\cdot625^2}{625^2\cdot961\cdot289}=\frac{30240}{961}$ Which our answer is $\fbox{961}$

 
                    \begin{figure}[H]
                        \centering
                        \includegraphics[width=0.5\textwidth]{Solution/AIME_2013_I_13_Solution_Figure_1.png}
                    \end{figure}
                    

 





\textbf{Solution 4}


                    \begin{figure}[H]
                        \centering
                        \includegraphics[width=0.5\textwidth]{Solution/AIME_2013_I_13_Solution_Figure_1.png}
                    \end{figure}
                    Let $k$ be the coefficient of the similarity of triangles 
 \[\triangle B_0 C_1 C_0 \sim \triangle AB_0 C_0 \implies k = \frac {B_0 C_0}{AC_0} = \frac {17}{25}.\]Then area $\frac {[B_0 C_1 C_0]}{[AB_0 C_0 ]} = k^2 \implies \frac {[AB_0 C_1]}{[AB_0 C_0]} = 1 – k^2.$ 

 The height of triangles $\triangle B_0C_1A$ and $\triangle AB_0C_0$ from $B_0$ is the same $\implies \frac {AC_1}{AC_0} = 1 – k^2.$

 The coefficient of the similarity of triangles $\triangle AB_1C_1 \sim \triangle AB_0C_0$ is  $\frac {AC_1}{AC_0} = 1 – k^2 \implies \frac {[B_1C_1C_2 ]}{[AB_0C_0 ]} = k^2 (1 – k^2)^2.$

 Analogically the coefficient of the similarity of triangles $\triangle AB_2C_2 \sim \triangle AB_0C_0$ is  $(1 – k^2)^2 \implies \frac {[B_2C_2C_3]}{[AB_0C_0 ]} = k^2 (1 – k^2)^4$ and so on.

 The yellow area $[Y]$ is $\frac {[Y]}{[AB_0C_0 ]} = k^2 + k^2 (1 – k^2)^2 + k^2 (1 – k^2)^4 +.. = \frac {k^2}{1 – (1 – k^2)^2} = \frac{1}{2 – k^2}.$

 The required area is $[AB_0C_0 ] – [Y] = [AB_0C_0 ] \cdot (1 –  \frac{1}{2 – k^2}) = [AB_0C_0 ] \cdot \frac  {1 –  k^2}{2 – k^2} = [AB_0C_0 ] \cdot \frac {25^2 – 17^2} {2 \cdot 25^2 – 17^2} =  [AB_0C_0 ] \cdot  \frac {336}{961}.$

 The number $961$ is prime, $[AB_0C_0]$ is integer but not $961,$ therefore the answer is $\boxed{961}$.

 \textbf{vladimir.shelomovskii@gmail.com, vvsss}

 


 


 




	\item (\textit{AIME I 2013}) For $\pi \le \theta < 2\pi$, let


 
 \[P=\dfrac12\cos\theta-\dfrac14\sin2\theta-\dfrac18\cos3\theta+\dfrac1{16}\sin4\theta+\dfrac1{32}\cos5\theta-\dfrac1{64}\sin6\theta-\dfrac1{128}\cos7\theta+\ldots\]


 and


 
 \[Q=1-\dfrac12\sin\theta-\dfrac14\cos2\theta+\dfrac1{8}\sin3\theta+\dfrac1{16}\cos4\theta-\dfrac1{32}\sin5\theta-\dfrac1{64}\cos6\theta+\dfrac1{128}\sin7\theta +\ldots\]


 so that $\frac{P}{Q} = \frac{2\sqrt2}{7}$. Then $\sin\theta = -\frac{m}{n}$ where $m$ and $n$ are relatively prime positive integers. Find $m+n$.


 
 \textbf{Answer: }036



\textbf{Solution 1}

Noticing the $\sin$ and $\cos$ in both $P$ and $Q,$ we think of the angle addition identities: 

 
 \[\sin(a + b) = \sin a \cos b + \cos a \sin b, \cos(a + b) = \cos a \cos b - \sin a \sin b\]

 With this in mind, we multiply $P$ by $\sin \theta$ and $Q$ by $\cos \theta$ to try and use some angle addition identities. Indeed, we get 
 \begin{align*} P \sin \theta + Q \cos \theta &= \cos \theta + \dfrac{1}{2}(\cos \theta \sin \theta - \sin \theta \cos \theta) - \dfrac{1}{4}(\sin{2 \theta} \sin \theta + \cos{2 \theta} \cos{\theta}) - \cdots \\ &= \cos \theta - \dfrac{1}{4} \cos \theta + \dfrac{1}{8} \sin{2 \theta} + \dfrac{1}{16} \cos{3 \theta} + \cdots \\ &= \cos \theta - \dfrac{1}{2}P \end{align*}
after adding term-by-term. Similar term-by-term adding yields 
 \[P \cos \theta + Q \sin \theta = -2(Q - 1).\] This is a system of equations; rearrange and rewrite to get 
 \[P(1 + 2 \sin \theta) + 2Q \cos \theta = 2 \cos \theta\] and 
 \[P \cos^2 \theta + Q \cos \theta(2 + \sin \theta) = 2 \cos \theta.\] Subtract the two and rearrange to get 
 \[\dfrac{P}{Q} = \dfrac{\cos \theta}{2 + \sin \theta} = \dfrac{2 \sqrt{2}}{7}.\] Then, square both sides and use Pythagorean Identity to get a quadratic in $\sin \theta.$ Factor that quadratic and solve for $\sin \theta = -17/19, 1/3.$ Since we're given $\pi\leq\theta<2\pi,$ $\sin\theta$ is nonpositive. We therefore use the negative solution, and our desired answer is $17 + 19 = \boxed{036}.$

 





\textbf{Solution 2}

Use sum to product formulas to rewrite $P$ and $Q$

 

 \[P \sin\theta\ + Q \cos\theta\ = \cos \theta\ - \frac{1}{4}\cos \theta + \frac{1}{8}\sin 2\theta + \frac{1}{16}\cos 3\theta - \frac{1}{32}\sin 4\theta + ...\]

 Therefore, 
 \[P \sin \theta - Q \cos \theta = -2P\]

 Using 
 \[\frac{P}{Q} =  \frac{2\sqrt2}{7}, Q = \frac{7}{2\sqrt2} P\]

 Plug in to the previous equation and cancel out the "P" terms to get: 
 \[\sin\theta - \frac{7}{2\sqrt2} \cos\theta = -2\]

 Then use the pythagorean identity to solve for $\sin\theta$, 
 \[\sin\theta = -\frac{17}{19} \implies \boxed{036}\]

 





\textbf{Solution 3}

Note that 
 \[e^{i\theta}=\cos(\theta)+i\sin(\theta)\]

 Thus, the following identities follow immediately:
 \[ie^{i\theta}=i(\cos(\theta)+i\sin(\theta))=-\sin(\theta)+i\cos(\theta)\]
 \[i^2 e^{i\theta}=-e^{i\theta}=-\cos(\theta)-i\sin(\theta)\]
 \[i^3 e^{i\theta}=\sin(\theta)-i\cos(\theta)\]

 Consider, now, the sum $Q+iP$. It follows fairly immediately that:

 
 \[Q+iP=1+\left(\frac{i}{2}\right)^1e^{i\theta}+\left(\frac{i}{2}\right)^2e^{2i\theta}+\ldots=\frac{1}{1-\frac{i}{2}e^{i\theta}}=\frac{2}{2-ie^{i\theta}}\]
 \[Q+iP=\frac{2}{2-ie^{i\theta}}=\frac{2}{2-(-\sin(\theta)+i\cos(\theta))}=\frac{2}{(2+\sin(\theta))-i\cos(\theta)}\]

 This follows straight from the geometric series formula and simple simplification. We can now multiply the denominator by it's complex conjugate to find:

 
 \[Q+iP=\frac{2}{(2+\sin(\theta))-i\cos(\theta)}\left(\frac{(2+\sin(\theta))+i\cos(\theta)}{(2+\sin(\theta))+i\cos(\theta)}\right)\]
 \[Q+iP=\frac{2((2+\sin(\theta))+i\cos(\theta))}{(2+\sin(\theta))^2+\cos^2(\theta)}\]

 Comparing real and imaginary parts, we find:
 \[\frac{P}{Q}=\frac{\cos(\theta)}{2+\sin(\theta)}=\frac{2\sqrt{2}}{7}\]

 Squaring this equation and letting $\sin^2(\theta)=x$:

 $\frac{P^2}{Q^2}=\frac{\cos^2(\theta)}{4+4\sin(\theta)+\sin^2(\theta)}=\frac{1-x^2}{4+4x+x^2}=\frac{8}{49}$

 Clearing denominators and solving for $x$ gives sine as $x=-\frac{17}{19}$.

 $017+019=\boxed{036}$

 





\textbf{Solution 4}

A bit similar to Solution 3. We use $\phi = \theta+90^\circ$ because the progression cycles in $P\in (\sin 0\theta,\cos 1\theta,-\sin 2\theta,-\cos 3\theta\dots)$. So we could rewrite that as $P\in(\sin 0\phi,\sin 1\phi,\sin 2\phi,\sin 3\phi\dots)$.

 Similarly, $Q\in (\cos 0\theta,-\sin 1\theta,-\cos 2\theta,\sin 3\theta\dots)\implies Q\in(\cos 0\phi,\cos 1\phi, \cos 2\phi, \cos 3\phi\dots)$.

 Setting complex $z=q_1+p_1i$, we get $z=\frac{1}{2}\left(\cos\phi+\sin\phi i\right)$

 $(Q,P)=\sum_{n=0}^\infty z^n=\frac{1}{1-z}=\frac{1}{1-\frac{1}{2}\cos\phi-\frac{i}{2}\sin\phi}=\frac{1-0.5\cos\phi+0.5i\sin\phi}{\dots}$.

 The important part is the ratio of the imaginary part $i$ to the real part. To cancel out the imaginary part from the denominator, we must add $0.5i\sin\phi$ to the numerator to make the denominator a difference (or rather a sum) of squares. The denominator does not matter. Only the numerator, because we are trying to find $\frac{P}{Q}=\tan\text{arg}(\Sigma)$ a PROPORTION of values. So denominators would cancel out.

 $\frac{P}{Q}=\frac{\text{Re}\Sigma}{\text{Im}\Sigma}=\frac{0.5\sin\phi}{1-0.5\cos\phi}=\frac{\sin\phi}{2-\cos\phi}=\frac{2\sqrt{2}}{7}$.

 Setting $\sin\theta=y$, we obtain
 \[\frac{\sqrt{1-y^2}}{2+y}=\frac{2\sqrt{2}}{7}\]
 \[7\sqrt{1-y^2}=2\sqrt{2}(2+y)\]
 \[49-49y^2=8y^2+32y+32\]
 \[57y^2+32y-17=0\rightarrow y=\frac{-32\pm\sqrt{1024+4\cdot 969}}{114}\]
 \[y=\frac{-32\pm\sqrt{4900}}{114}=\frac{-16\pm 35}{57}\].

 Since $y<0$ because $\pi<\theta<2\pi$, $y=\sin\theta=-\frac{51}{57}=-\frac{17}{19}$. Adding up, $17+19=\boxed{036}$.

 





\textbf{Solution 5 (utterly disgusting)}

We notice $\sin\theta=-\frac{i}{2}(e^{i\theta}-e^{-i\theta})$ and $\cos\theta=\frac{1}{2}(e^{i\theta}+e^{-i\theta})$

 We observe that both $P$ and $Q$ can be split into $2$ parts, namely the terms which contain the $\cos$ and the terms which contain the $\sin .$

 The $\cos$ part of $P$ can be expressed as: 
 \begin{align*}\frac12\cos\theta-\frac18\cos3\theta+\cdots&=\frac14\left(e^{i\theta}\left(1-\frac{e^{i2\theta}}{4}+\cdots\right)+e^{-i\theta}\left(1-\frac{e^{-i2\theta}}{4}+\cdots\right)\right) \\ &= \frac{1}{4}\left(\frac{e^{i\theta}}{1+\frac{1}{4}e^{i2\theta}}+\frac{e^{-i\theta}}{1+\frac{1}{4}e^{-i2\theta}}\right)\\  &= \frac{5(e^{i\theta}+e^{-i\theta})}{17+4e^{i2\theta}+4e^{-i2\theta}}.\end{align*}


 Repeating the above process, we find that the $\sin$ part of $P$ is 
 \[\frac{2i(e^{i2\theta}-e^{-i2\theta})}{17+4e^{i2\theta}+4e^{-i2\theta}},\] the $\cos$ part of $Q$ is 
 \[\frac{16+2(e^{i2\theta}+e^{-i2\theta})}{17+4e^{i2\theta}+4e^{-i2\theta}},\] and finally, the $\sin$ part of $Q$ is 
 \[\frac{3i(e^{i\theta}-e^{-i\theta})}{17+4e^{i2\theta}+4e^{-i2\theta}}.\]

 Converting back to trigonometric form, we have 
 \begin{align*}\frac{2\sqrt{2}}{7}&=\frac{10\cos{\theta}-4\sin{2\theta}}{16+4\cos{2\theta}-6\sin{\theta}}\\ &=\frac{5\cos{\theta}-2\sin{2\theta}}{8+2\cos{2\theta}-3\sin{\theta}}.\end{align*}
 Using the $\sin$ double identity and simplifying, we have 
 \[\frac{2\sqrt2}{7}=\frac{\cos{\theta}(5-4\sin{\theta})}{10-4\sin^2{\theta}-3\sin{\theta}}.\] Factoring the denominator, we have 
 \[10-4\sin^2{\theta}-3\sin{\theta}=(5-4\sin\theta)(2+\sin\theta).\] Simplifying 
 \begin{align*}\frac{2\sqrt2}{7}&= \frac{\cos{\theta}(5-4\sin{\theta})}{(5-4\sin\theta)(2+\sin\theta)}\\ &=\frac{\cos\theta}{2+\sin\theta}.\end{align*}
 We set $\sin\theta$ as $x$, and by the Pythagorean Identity, we have $57x^2+32x-17=0$. This factors into $(19x+17)(3x-1)=0$, which yields the 2 solutions $x=-\frac{17}{19}, x=\frac{1}{3}$. As $\pi\leq\theta<2\pi$, the latter root is erroneous, and we are left with $\sin\theta=-\frac{17}{19}$. Thus, our final answer is $17+19=\boxed{036}$.

 





\textbf{Solution 6}

Follow solution 3, up to the point of using the geometric series formula
 \[Q+iP=\frac{1}{1+\frac{\sin(\theta)}{2}-\frac{Qi\cos(\theta)}{2}}\]

 Moving everything to the other side, and considering only the imaginary part, we get
 \[Pi+\frac{Pi}{2}\sin\theta-\frac{Qi}{2}\cos\theta = 0\]

 We can then write $P = 2 \sqrt{2} k$, and $Q = 7k$, ($k \neq 0$). Thus, we can substitute and divide out by k.
 \[2\sqrt{2}+\sqrt{2}\sin\theta-\frac{7}{2}\cos\theta\ =\ 0\]
 \[2\sqrt{2}+\sqrt{2}\sin\theta-\frac{7}{2}\sqrt{1-\sin^{2}\theta}=\ 0\]
 \[2\sqrt{2}+\sqrt{2}\sin\theta\ =\frac{7}{2}\left(\sqrt{1-\sin^{2}\theta}\right)\]
 \[8+8\sin\theta+2\sin^{2}\theta=\frac{49}{4}-\frac{49}{7}\sin^{2}\theta\]
 \[\frac{57}{4}\sin^{2}\theta+8\sin\theta-\frac{17}{4} = 0\]
 \[57\sin^{2}\theta+32\sin\theta-17 = 0\]
 \[\left(3\sin\theta-1\right)\left(19\sin\theta+17\right) = 0\]

 Since $\pi \le \theta < 2\pi$, we get $\sin \theta < 0$, and thus, $\sin\theta = \frac{-17}{19} \implies \boxed{036}$

 -Alexlikemath

 


 




	\item (\textit{AIME I 2013}) Let $N$ be the number of ordered triples $(A,B,C)$ of integers satisfying the conditions: 


 (a) $0\le A<B<C\le99$, 


 (b) there exist integers $a$, $b$, and $c$, and prime $p$ where $0\le b<a<c<p$, 


 (c) $p$ divides $A-a$, $B-b$, and $C-c$, and 


 (d) each ordered triple $(A,B,C)$ and each ordered triple $(b,a,c)$ form arithmetic sequences. Find $N$.


 
 \textbf{Answer: }272



\textbf{Solution}

From condition (d), we have $(A,B,C)=(B-D,B,B+D)$ and $(b,a,c)=(a-d,a,a+d)$.

 
Now, we substitute these values in. Condition $\text{(c)}$ states that $p\mid B-D-a$, $p | B-a+d$, and $p\mid B+D-a-d$. We subtract the first two to get $p\mid-d-D$, and we do the same for the last two to get $p\mid 2d-D$. We subtract these two to get $p\mid 3d$. So $p\mid 3$ or $p\mid d$. The second case is clearly impossible, because that would make $c=a+d>p$, violating condition $\text{(b)}$. So we have $p\mid 3$, meaning $p=3$. Condition $\text{(b)}$ implies that $(b,a,c)=(0,1,2)$ or $(a,b,c)\in (1,0,2)\rightarrow (-2,0,2)\text{ }(D\equiv 2\text{ mod 3})$. 

 
Now we return to condition $\text{(c)}$, which now implies that $(A,B,C)\equiv(-2,0,2)\pmod{3}$. Now, we set $B=3k$ for increasing positive integer values of $k$. $B=0$ yields no solutions. $B=3$ gives $(A,B,C)=(1,3,5)$, giving us $1$ solution. If $B=6$, we get $2$ solutions, $(4,6,8)$ and $(1,6,11)$. Proceeding in the manner, we see that if $B=48$, we get 16 solutions. However, $B=51$ still gives $16$ solutions because $C_\text{max}=2B-1=101>100$. Likewise, $B=54$ gives $15$ solutions. This continues until $B=96$ gives one solution. $B=99$ gives no solution. Thus, $N=1+2+\cdots+16+16+15+\cdots+1=2\cdot\frac{16(17)}{2}=16\cdot 17=\boxed{272}$.

 





\textbf{Solution 2}

Let $(A, B, C) = (B-x, B, B+x)$ and $(b, a, c) = (a-y, a, a+y)$. Now the 3 differences would be 
 \begin{align} \label{1} &A-a = B-x-a \\ \label{2} &B - b = B-a+y \\ \label{3} &C - c = B+x-a-y  \end{align}

 Adding equations $(1)$ and $(3)$ would give $2B - 2a - y$. Then doubling equation $(2)$ would give $2B - 2a + 2y$. The difference between them would be $3y$. Since $p|\{(1), (2), (3)\}$, then $p|3y$. Since $p$ is prime, $p|3$ or $p|y$. However, since $p > y$, we must have $p|3$, which means $p=3$.

 
If $p=3$, the only possible values of $(b, a, c)$ are $(0, 1, 2)$. Plugging this into our differences, we get 
 \begin{align*}  &A-a = B-x-1 \hspace{4cm}(4)\\  &B - b = B \hspace{5.35cm}(5)\\  &C - c = B+x-2 \hspace{4cm}(6) \end{align*}
The difference between $(4)$ and $(5)$ is $x+1$, which should be divisible by 3. So $x \equiv 2 \mod 3$. Also note that since $3|(5)$, $3|B$. Now we can try different values of $x$ and $B$:

 When $x=2$, $B=3, 6, ..., 96 \Rightarrow 32$ triples.

 When $x=5$, $B=6, 9, ..., 93\Rightarrow 30$ triples..

 ... and so on until

 When $x=47$, $B=48, 51\Rightarrow 2$ triple. 

 So the answer is $32 + 30 + \cdots + 2 = \boxed{272}$

 





\textbf{Solution 3}

Let $d=B-A=C-B$ and $e=a-b=c-a$ be the common difference of the two sequences. We have \begin{align*}A=a+pi\\B=b+pj\\C=c+pk\end{align*}
For some integers $i,j,k$. Computing $B-A$ and $C-B$, we have\begin{align*}B-A=d=(b-a)+p(j-i)\\C-B=d=(c-b)+p(k-j)\end{align*}
Subtracting the two equations again, we find
 \[(a+c-2b)+p(i+k-2j)=0\]Or
 \[p(2j-i-k)=3e\]Clearly $0<e<p$. It follows that $p\mid 3e\implies p=3$. 

 We have $(b,a,c)=(0,1,2)$, which gives 
 \[A\equiv 1\bmod{3},\quad B\equiv 0\bmod{3},\quad C\equiv 2\bmod{3}\]Thus, $2d\equiv 1\bmod{3}$. Let $2d=1+3x$, which needs to be even. Substituting $x=2y-1$, we have $2d=6y-2$. 

 An order-triple $(A,B,C)$ is uniquely determined by a solution of $A+2d\le 99, A\equiv 1\bmod{3}$, while satisfying condition (a). It follows that
 \[A+6y-2\le 99\iff y\le\frac{101-A}{6}\]The number of solutions for indices $y$ given an $A$-value is
 \[\left\lfloor \frac{101-A}{6}\right\rfloor=16+\left\lfloor\frac{5-A}{6}\right\rfloor\]For $A\in \{1,4,\ldots, 94\}$, the number of order-triples $(A,B,C)$ is
 \[16+16+15+15+\ldots+1+1=\boxed{272}\]

 




\end{enumerate}